%------------------------------------------------------------------------------%
\chapter{Matrizes}
\label{cap:matrizes}

%------------------------------------------------------------------------------%
\section{Matrizes e operações matriciais }

%------------------------------------------------------------------------------%
\begin{teoria}{Matrizes}
    
    Definimos como matriz do tipo $m\times n$ o conjunto de números reais, dispostos em um quadro de $m$ linhas por $n$ colunas.
    \[ A = \left[
    \begin{array}{ccccc}
        a_{11} & a_{12} & a_{13} & \cdots & a_{1n} \\
        a_{21} & a_{22} & a_{23} & \cdots & a_{2n} \\
        a_{31} & a_{32} & a_{33} & \cdots & a_{3n} \\
        \vdots & \vdots & \vdots & \vdots & \vdots \\
        a_{m1} & a_{m2} & a_{m3} & \cdots & a_{mn} 
    \end{array}
    \right] \]
    Podemos definir uma matriz abreviadamente como $A=(a_{ij})_{m\times n}$.
    Os valores $a_{ij}$ são chamados \textbf{elementos} de $A$.
    \\
    
    Dizemos que uma Matriz é \textbf{Quadrada} se $m=n$. 
    \\
    
    Uma matriz quadrada importante é a \textbf{Matriz Identidade}. Essas matrizes tem todos os elementos da diagonal iguais a $1$ e os demais elementos nulos
    \[ 
        I_2 = \left[\begin{array}{cc}
            1 & 0 \\ 
            0 & 1
        \end{array}\right] 
    \qquad\qquad
         I_3 = \left[\begin{array}{ccc}
             1 & 0 & 0 \\ 
             0 & 1 & 0 \\
             0 & 0 & 1 
         \end{array}\right] 
     \]
     
     Dizemos que duas matrizes são iguais se elas tem as mesmas dimensões e todos os seus elementos são iguais.
     
\end{teoria}

%------------------------------------------------------------------------------%
\begin{teoria}{Transposição}
    
    \textbf{Matriz Transposta} ou transposta de uma matriz $A=(a_{ij})_{m\times n}$ é uma matriz construída trocando as linhas e colunas de~$A$
    \[ A^t = (a_{ji})_{n\times m} \]
    
\end{teoria}

%------------------------------------------------------------------------------%
\begin{teoria}{Soma}
    
    A \textbf{Soma} de duas matrizes de mesma dimensão $A=(a_{ij})_{m\times n}$ e $B=(b_{ij})_{m\times n}$
    é uma matriz construída somando cada elemento de $A$ com o elemento correspondente de $B$. 
    Assim \[C=A+B\] significa que
    \[ C = (c_{ij})_{m\times n} = (a_{ij}+b_{ij})_{m\times n} =
        \left[
        \begin{array}{cccc}
            a_{11} + b_{11} & a_{12} + b_{12} & \cdots & a_{1n} + b_{1n} \\
            a_{21} + b_{21} & a_{22} + b_{22} & \cdots & a_{2n} + b_{2n} \\
            \vdots          & \vdots          & \vdots & \vdots          \\
            a_{m1} + b_{m1} & a_{m2} + b_{m2} & \cdots & a_{mn} + b_{mn}
        \end{array}
        \right]
    \]
        
\end{teoria}

%------------------------------------------------------------------------------%
\begin{teoria}{Subtração}
    
    A \textbf{Subtração} de matrizes é definida de modo equivalente a soma, ou seja, \[C=A-B\] significa que
    \[ C = (c_{ij})_{m\times n} = (a_{ij}-b_{ij})_{m\times n} =
        \left[
        \begin{array}{cccc}
            a_{11} - b_{11} & a_{12} - b_{12} & \cdots & a_{1n} - b_{1n} \\
            a_{21} - b_{21} & a_{22} - b_{22} & \cdots & a_{2n} - b_{2n} \\
            \vdots          & \vdots          & \vdots & \vdots          \\
            a_{m1} - b_{m1} & a_{m2} - b_{m2} & \cdots & a_{mn} - b_{mn}
        \end{array}
        \right]
    \]

\end{teoria}

%------------------------------------------------------------------------------%
\begin{teoria}{Multiplicação por um número real}
        
    \textbf{Multiplicação por um número real} ou Multiplicação por um Escalar é uma operação que multiplica cada elemento da matriz por um número, 
    isso é, se $k$ é um número real, então \[B=kA\] é a operação
    \[ B = (b_{ij})_{m\times n} = (ka_{ij})_{m\times n} =
        \left[
        \begin{array}{cccc}
            k a_{11} & k a_{12} & \cdots & k a_{1n} \\
            k a_{21} & k a_{22} & \cdots & k a_{2n} \\
            \vdots   &   \vdots & \vdots &   \vdots \\
            k a_{m1} & k a_{m2} & \cdots & k a_{mn}
        \end{array}
        \right]
    \]    
    
\end{teoria}

%------------------------------------------------------------------------------%
\begin{teoria}{Multiplicação de Matrizes}
    
    A \textbf{Multiplicação de Matrizes} é uma operação menos intuitiva do que as demais. 
    Ela é definida entre duas matrizes tais que o número de linhas da segunda seja igual ao número de colunas da primeira.
    Assim para multiplicarmos as matrizes $A$ e $B$ elas devem ter as dimensões compatíveis
    \[ A=(a_{ij})_{m\times n} \qquad\qquad B=(b_{ij})_{n\times p} \]
    A matriz \[C=AB\] terá dimensões $C = (c_{ij})_{m\times p}$ e seu elemento $c_{ij}$ é igual a soma dos produtos da $i$-ésima linha de $A$ 
    com os elementos da $j$-ésima coluna de $B$.
    \[ c_{ij} = \sum_{k=1}^{n} a_{ik}\,b_{kj}  \]
    Por exemplo, se $A$ for uma matriz $2\times3$ e $B$ uma matriz $3\times2$ teremos
    \[
        \left[\begin{array}{ccc}
            a & b & c \\ 
            d & e & f
        \end{array} \right]
        \times
        \left[\begin{array}{cc}
            t & u \\ 
            v & x \\ 
            y & z
        \end{array} \right]
        =
        \left[\begin{array}{cc}
            at+bv+cy & au+bx+cz \\ 
            dt+ev+fy & du+ex+fz
        \end{array} \right]
    \]
    
    A multiplicação de uma matriz pela matriz identidade retorna a matriz original
    \[ AI_n = A \qquad\qquad I_mA=A  \]
    
\end{teoria}

%------------------------------------------------------------------------------%
\section{Propriedades das Operações com Matrizes}

%------------------------------------------------------------------------------%
\begin{teoria}{Propriedades das Operações com Matrizes -- Soma}
    
    Sejam $A$, $B$ e $C$ matrizes com as dimensões apropriadas, 
    então as operações com matrizes tem as seguintes propriedades
    \begin{propriedades}
        \item $A+B = B+A$                              \tabto{8cm} Comutatividade da soma
        \item $A+(B+C) = (A+B)+C$                      \tabto{8cm} Associatividade da soma
        \item Se $\mathbf{0}$ é uma matriz com todos os elementos nulos \\ então
              $A+\mathbf{0}=A$                         \tabto{8cm} Elemento neutro da soma
        \item Para cada matriz $A$ existe uma matriz \\ $-A=(-1)A$ tal que
              $A+(-A)=\mathbf{0}$                      \tabto{8cm} Elemento simétrico
    \end{propriedades}
    
\end{teoria}

%------------------------------------------------------------------------------%
\begin{teoria}{Propriedades das Operações com Matrizes -- Produto por Escalar}
    
    Sejam $A$, $B$ e $C$ matrizes com as dimensões apropriadas e $\alpha$ e $\beta$ números reais (escalares), 
    então as operações com matrizes tem as seguintes propriedades
    \begin{propriedades}
        \item $\alpha(\beta A)=(\alpha\beta)A$         \tabto{8cm} Associatividade do produto por escalar
        \item $(\alpha+\beta)A= \alpha A + \beta A$    \tabto{8cm} Distributividade do produto por escalar
        \item $\alpha(A+B)=\alpha A+\alpha B$          \tabto{8cm} Distributividade do produto por escalar
    \end{propriedades}
    
\end{teoria}

%------------------------------------------------------------------------------%
\begin{teoria}{Propriedades das Operações com Matrizes -- Produto}
    
    Sejam $A$, $B$ e $C$ matrizes com as dimensões apropriadas e $\alpha$ e $\beta$ números reais (escalares), 
    então as operações com matrizes tem as seguintes propriedades
    \begin{propriedades}
        \item $A(BC) = (AB)C$                          \tabto{8cm} Associatividade do produto
        \item $A(B+C) = AB+AC$                         \tabto{8cm} Distributividade do produto
        \item $(B+C)A = BA+CA$                         \tabto{8cm} Distributividade do produto
        \item A matriz identidade $I_n$ é tal que \\
        $AI_n=I_mA=A$                            \tabto{8cm} Elemento neutro do produto
        \item $\alpha(AB) = (\alpha A)B = A(\alpha B)$ \tabto{8cm} Associatividade dos produtos
    \end{propriedades}
    
\end{teoria}

\begin{observacao}
    
    A multiplicação de matrizes não é comutativa, isso é, em geral
    \[ AB \neq BA \]
    Na multiplicação de matrizes não vale a lei do anulamento do produto, é possível que $A\neq0$ e $B\neq0$ e mesmo assim $AB=0$
    \\
    
    Se acontecer $AB=BA$ dizemos que as matrizes $A$ e $B$ são comutativas.
    
\end{observacao}


%------------------------------------------------------------------------------%
\begin{teoria}{Propriedades das Operações com Matrizes -- Transposição}
    
    Sejam $A$, $B$ e $C$ matrizes com as dimensões apropriadas e $\alpha$ e $\beta$ números reais (escalares), 
    então as operações com matrizes tem as seguintes propriedades
    \begin{propriedades}
        \item $(A^t)^t=A$
        \item $(A+B)^t =A^t + B^t$
        \item $(\alpha A)^t = \alpha A^t$
        \item $(AB)^t=B^tA^t$
    \end{propriedades}
    
\end{teoria}

%------------------------------------------------------------------------------%
\begin{teoria}{Simetria}

    Dizemos que uma matriz $A$ é 
    \begin{description}
        \item[Simétrica]      \tabto{20mm} quando $A^t=A$
        \item[Antissimétrica] \tabto{20mm} quando $A^t=-A$
    \end{description}
    
\end{teoria}

%------------------------------------------------------------------------------%
\begin{exemplos}
    
    \begin{exemplo}
        Verifique se a identidade é verdadeira ou falsa
        \[ (A+B)(A-B) = A^2 - B^2 \]
        
        Resolução:    
        \begin{align*}
            (A+B)(A-B) & = (A+B)A-(A+B)B \\
                       & = AA + BA - AB - BB \\
                       & = A^2 + BA - AB - B^2
        \end{align*}
        A identidade só será verdadeira se $BA - AB=0$ o que não é verdade em geral para matrizes. 
        Portanto a identidade não é verdadeira.
        
        Fazendo as contas com um exemplo numérico simples
        \[ A = \left[\begin{array}{cc}
            0 & 0 \\ 
            1 & 1
            \end{array} \right] 
            \qquad
            B = \left[\begin{array}{cc}
            1 & 0 \\ 
            1 & 0
            \end{array} \right] 
        \]
        Calculamos agora $(A+B)(A-B)$ e $A^2 - B^2$ e verificamos que os resultados não serão iguais
        \[ A+B = \left[\begin{array}{cc}
                0 & 0 \\ 
                1 & 1
                \end{array} \right] 
                +
                \left[\begin{array}{cc}
                1 & 0 \\ 
                1 & 0
                \end{array} \right]
              = \left[\begin{array}{cc}
                1 & 0 \\ 
                2 & 1
                \end{array} \right]  
        \]
        \[ A-B = \left[\begin{array}{cc}
                0 & 0 \\ 
                1 & 1
                \end{array} \right] 
                -
                \left[\begin{array}{cc}
                1 & 0 \\ 
                1 & 0
                \end{array} \right]
              = \left[\begin{array}{rc}
                -1 & 0 \\ 
                 0 & 1
                \end{array} \right]  
        \]    
        \[ (A+B)(A-B) = \left[\begin{array}{cc}
                1 & 0 \\ 
                2 & 1
                \end{array} \right]
                \,
                \left[\begin{array}{rc}
                -1 & 0 \\ 
                0 & 1
                \end{array} \right]
              = \left[\begin{array}{rc}
                -1 & 0 \\ 
                -2 & 1
                \end{array} \right]  
        \]
        \[ A^2 = \left[\begin{array}{cc}
                1 & 0 \\ 
                1 & 0
                \end{array} \right]^2
              = \left[\begin{array}{cc}
                1 & 0 \\ 
                1 & 0
                \end{array} \right]
                \,
                \left[\begin{array}{cc}
                1 & 0 \\ 
                1 & 0
                \end{array} \right]
              = \left[\begin{array}{cc}
                0 & 0 \\ 
                1 & 1
                \end{array} \right]  
        \]
        \[ B^2 = \left[\begin{array}{cc}
                0 & 0 \\ 
                1 & 1
                \end{array} \right]^2
              = \left[\begin{array}{cc}
                0 & 0 \\ 
                1 & 1
                \end{array} \right]
                \,
                \left[\begin{array}{cc}
                0 & 0 \\ 
                1 & 1
                \end{array} \right]
              = \left[\begin{array}{cc}
                1 & 0 \\ 
                1 & 0
                \end{array} \right]  
        \]
        \[ A^2-B^2 = \left[\begin{array}{cc}
                0 & 0 \\ 
                1 & 1
                \end{array} \right]
                -
                \left[\begin{array}{cc}
                1 & 0 \\ 
                1 & 0
                \end{array} \right]
                = \left[\begin{array}{rc}
                -1 & 0 \\ 
                 0 & 1
                \end{array} \right]  
        \]
    \end{exemplo}
                    
\end{exemplos}

%------------------------------------------------------------------------------%
\begin{teoria}{Matriz Inversa}
    
    Se $A$ é uma matriz quadrada de ordem $n$ se existir uma matriz $B$ com a mesma ordem que $A$ e que satisfaça
    \[ AB = BA = I_n \]
    a matriz $B$ é chamada de \textbf{Inversa} da matriz $A$.    
    Se a inversa da matriz $A$ existe podemos indicá-la por $A^{-1}$. 
    \\
    
    Quando $A$ não possui uma inversa dizermos que $A$ é uma \textbf{Matriz Singular}.
    
\end{teoria}

%------------------------------------------------------------------------------%
\section{Determinante de uma Matriz e suas Propriedades}

%------------------------------------------------------------------------------%
\begin{teoria}{Determinante de uma Matriz Quadrada de Ordem 1}

  O determinante de uma matriz com uma linha e uma coluna é igual ao valor do elemento da matriz
  \[  A = [a_{11}]  \qquad \text{ então } \qquad  \det A = a_{11} \]

\end{teoria}

%------------------------------------------------------------------------------%
\begin{teoria}{Determinante de uma Matriz Quadrada de Ordem 2}
    
    Data uma matriz $A$ quadrada de ordem 2 
    \[ A = \left[
        \begin{array}{cc}
            a & b \\ 
            c & d
        \end{array} 
        \right]
    \]
    definimos o determinante de $A$ como sendo a subtração do produto dos elementos da diagonal principal menos o produto dos elementos da diagonal secundária
    \[ \det A = \left|
        \begin{array}{cc}
            a & b \\ 
            c & d
        \end{array} 
        \right|
        = ad - bc    
    \]
    
    Note a diferença na notação de Matriz $[\;\cdot\;]$ e de determinante $|\,\cdot\,|$.
    
\end{teoria}

%------------------------------------------------------------------------------%
\begin{teoria}{Determinante de uma Matriz Quadrada de Ordem 3}
    
    Podemos calcular o determinante de uma matriz quadrada de ordem 3 pela \textbf{Regra de Sarrus}.
    Considere a matriz
    \[ A = \left[\begin{array}{ccc}
        a_{11} & a_{12} & a_{13} \\ 
        a_{21} & a_{22} & a_{23} \\ 
        a_{31} & a_{32} & a_{33}
        \end{array}\right]
    \]
    Copiamos as duas primeiras colunas no final da matriz e 
    somamos os produtos dos elementos das diagonais à direita e subtraímos os produtos dos elementos das diagonais à esquerda

    \addgraphic{0.3}{regra-sarrus}

    Obtemos assim o resultado
    \[ \det A = a_{11} a_{22} a_{33} + a_{12} a_{23} a_{33} + a_{13} a_{21} a_{32} 
               -a_{13} a_{22} a_{31} - a_{11} a_{23} a_{32} - a_{12} a_{21} a_{33}
     \]
     
     \begin{observacao}
         Note que essa regra se aplica apenas a matrizes de ordem 3.
     \end{observacao}
    
\end{teoria}

\begin{observacao}
    Para definirmos o determinante para matrizes quadradas de ordem maior do que 3 precisamos de algumas definições auxiliares e um teorema.
\end{observacao}

%------------------------------------------------------------------------------%
\begin{teoria}{Menor Complementar}

    O \textbf{Menor complementar}, $\minor{A}_{ij}$, de uma matriz $A=(a_{ij})_{n\times n}$ relativo ao elemento $a_{ij}$ é o \textbf{determinante}
    da matriz obtida removendo a linha $i$ e a coluna $j$ da matriz $A$. 

\end{teoria}

\begin{exemplos}
    \begin{exemplo}
        Se $A$ for a matriz
        \[ A = \left[\begin{array}{ccc}
        a_{11} & a_{12} & a_{13} \\ 
        a_{21} & a_{22} & a_{23} \\ 
        a_{31} & a_{32} & a_{33}
        \end{array}\right]
        \]
        Então esses são alguns dos seus Menores Complementares
        \[  \minor{A}_{11} = \left|\begin{array}{cc}
        	a_{22} & a_{23} \\
        	a_{32} & a_{33}
        \end{array}\right| = a_{22} a_{33} - a_{23} a_{32} 
        \]
        \[ \minor{A}_{12} = \left|\begin{array}{cc}
        	a_{21} & a_{23} \\
        	a_{31} & a_{33}
        \end{array}\right|
        = a_{21} a_{33} - a_{23} a_{31}
        \]
        \[ \minor{A}_{23} = \left|\begin{array}{cc}
        	a_{11} & a_{12} \\
        	a_{31} & a_{32}
        \end{array}\right|
        = a_{11} a_{32} - a_{12} a_{31}
        \]
    \end{exemplo}
\end{exemplos}

%------------------------------------------------------------------------------%
\begin{teoria}{Cofator}

    O \textbf{Cofator} do elemento $a_{ij}$ é o número obtido multiplicando $(-1)^{i+j}$ com o Menor Complementar de $a_{ij}$
    \[ A_{ij} = (-1)^{i+j} \minor{A}_{ij} \]
    
\end{teoria}

%------------------------------------------------------------------------------%
\begin{teoria}{Determinante de uma Matriz Qualquer Ordem -- Teorema de Laplace}
    
    O \textbf{Teorema de Laplace} nos diz como calcular o determinante de uma matriz $A$. \\
    
    O determinante é calculado somando os Cofatores de cada elemento de uma linha ou coluna da matriz.\\
    
    Se aplicarmos o teorema a primeira linha temos
    \[ \det A = a_{11} A_{11} + a_{12} A_{12} + a_{13} A_{13} + \cdots + a_{1n} A_{1n} \]
    Explicitando os Menores Complementares temos
    \[ \det A = a_{11}(-1)^{1+1}\minor{A}_{11} + 
                a_{12}(-1)^{1+2}\minor{A}_{12} + 
                a_{13}(-1)^{1+3}\minor{A}_{13} + \cdots + 
                a_{1n}(-1)^{1+n}\minor{A}_{1n} \]

\end{teoria}

%------------------------------------------------------------------------------%
\begin{exemplos}
    
    \begin{exemplo}
        Calcule o determinante da matriz
        \[ M = \left[
            \begin{array}{cccc}
                4 & x & 0 & x \\ 
                3 & 0 & 2 & 0 \\ 
                0 & 2 & 0 & 0 \\ 
                x & 0 & 1 & 0
            \end{array} 
            \right]
        \]
        
        Resolução:
        \begin{align*}
        \det M & = \left|
                        \begin{array}{cccc}
                        4 & x & 0 & x \\ 
                        3 & 0 & 2 & 0 \\ 
                        0 & 2 & 0 & 0 \\ 
                        x & 0 & 1 & 0
                        \end{array} 
                        \right| \\
            & =  a_{11} A_{11} + a_{12} A_{12} + a_{13} A_{13} + a_{14} A_{14} \\
            & =  a_{11} \minor{A}_{11} - a_{12} \minor{A}_{12} + a_{13} \minor{A}_{13} - a_{14} \minor{A}_{14} \\
            & = 4 \underbrace{\left|
                    \begin{array}{ccc}
                    0 & 2 & 0 \\ 
                    2 & 0 & 0 \\ 
                    0 & 1 & 0
                    \end{array} 
                    \right|}_{=0}
              - x \underbrace{\left|
                    \begin{array}{ccc}
                    3 & 2 & 0 \\ 
                    0 & 0 & 0 \\ 
                    x & 1 & 0
                    \end{array} 
                    \right|}_{=0}
              + 0 \left|
                    \begin{array}{ccc}
                    3 & 0 & 0 \\ 
                    0 & 2 & 0 \\ 
                    x & 0 & 0
                    \end{array} 
                    \right|
              - x \left|
                    \begin{array}{ccc}
                    3 & 0 & 2 \\ 
                    0 & 2 & 0 \\ 
                    x & 0 & 1
                    \end{array} 
                    \right| 
              = - x \left|
                    \begin{array}{ccc}
                    3 & 0 & 2 \\ 
                    0 & 2 & 0 \\ 
                    x & 0 & 1
                    \end{array} 
                    \right| \\
            & = - x \left( \;
                      3 \left|
                        \begin{array}{cc}
                        2 & 0 \\ 
                        0 & 1
                        \end{array} 
                        \right|
                     -0 \left|
                        \begin{array}{cc}
                        0 & 0 \\ 
                        x & 1
                        \end{array} 
                        \right|
                     +2 \left|
                        \begin{array}{cc}
                        0 & 2 \\ 
                        x & 0
                        \end{array} 
                        \right| \;
                    \right) \\
            & = -x \big[ 3(2\times 1-0) +2(0-2x) \big] 
              = -x (6-4x) 
              = 2x(x-3)
        \end{align*}
    \end{exemplo}
\end{exemplos}

%------------------------------------------------------------------------------%
\begin{teoria}{Combinação Linear}
    
    Uma \textbf{Combinação Linear} de colunas de uma matriz é a soma de duas colunas podendo multiplicar cada uma delas por um número não negativo.
    Por exemplo, considerando a matriz $A$
    \[ A = \left[\begin{array}{ccc}
        a_{11} & a_{12} & a_{13} \\ 
        a_{21} & a_{22} & a_{23} \\ 
        a_{31} & a_{32} & a_{33}
        \end{array}\right]
    \]
    Podemos construir uma matriz $B$ fazendo uma combinação linear da primeira e segunda colunas de $A$
    \[ B = \left[\begin{array}{ccc}
        a_{11} & r a_{12} + s a_{11}& a_{13} \\ 
        a_{21} & r a_{22} + s a_{21}& a_{23} \\ 
        a_{31} & r a_{32} + s a_{31}& a_{33}
    \end{array}\right]
    \]
    onde $r$ e $s$ são números reais não nulos.
    O mesmo pode ser feito com as linhas de uma matriz.
    
\end{teoria}

%------------------------------------------------------------------------------%
\begin{teoria}{Propriedades dos Determinantes}
    
    O determinante de uma matriz será zero quando
    \begin{enumerate}
        \item Todos os elementos de uma de suas linhas (ou colunas) forem zero
        \item Quando a matriz possui duas linhas (ou colunas) iguais ou proporcionais
        \item Quando uma linha (ou coluna) é a combinação linear de outras linhas (ou colunas)
    \end{enumerate}
    
    Transformações que não alteram o determinante de uma matriz
    \begin{enumerate}
        \item Quando transpomos a matriz \[\det A^t = \det A\]
        \item Quando somamos aos elementos de uma linha (ou coluna) os elementos de outra linha (ou coluna) multiplicados por uma constante
    \end{enumerate}
    
    Transformações que alteram o determinante
    \begin{enumerate}
        \item O determinante muda de sinal se trocarmos duas linhas (ou colunas) de lugar
        \item Se multiplicarmos linha (ou coluna) por uma constante o valor do determinante será multiplicado por essa mesma constante.
    \end{enumerate}
    
    Se $A$ e $B$ forem duas matrizes quadradas de mesma ordem, então
    \[ \det(AB) = \det A \det B \]
    
    O determinante de qualquer matriz identidade é sempre $1$
    \[ \det I = 1 \]
    
    Uma matriz quadrada $A$ possui inversa se, e somente se, 
    \[\det A\neq 0\]
    
    O determinante da inversa da matriz $A$ é
    \[ \det A^{-1} = \frac{1}{\;\det A\;} \]
    
\end{teoria}

%------------------------------------------------------------------------------%
\section{Sistemas Lineares}

%------------------------------------------------------------------------------%
\begin{teoria}{Sistemas Lineares}
    Uma sistema linear  é um conjunto de $m$ equações lineares com $n$ incógnitas $x_1,\,x_2,\,\dots,\,x_n$
    \[\left\{ 
    \begin{array}{ccccccccl}
        a_{11}x_1 & \!+\! &a_{12}x_2  & \!+\! &\cdots & \!+\! &a_{1n}x_n & \!=\! & b_1 \\ 
        a_{21}x_1 & \!+\! &a_{22}x_2  & \!+\! &\cdots & \!+\! &a_{2n}x_n & \!=\! & b_2 \\ 
        \vdots    &       &\vdots     &       &       &       &\vdots    &       & \vdots \\
        a_{m1}x_1 & \!+\! &a_{m2}x_2  & \!+\! &\cdots & \!+\! &a_{mn}x_n & \!=\! & b_m
    \end{array} 
    \right.
    \]
\end{teoria}

\begin{teoria}{}    
    Em termos de número de soluções, um sistema linear pode ser classificado como
    \begin{description}
        \item[Possível] Admite pelo menos uma solução
        \begin{description}
            \item[Determinado]   \tabto{22mm} Admite uma única solução
            \item[Indeterminado] \tabto{22mm} Admite infinitas soluções
        \end{description}
        \item[Impossível] Não existe nenhuma solução para o sistema 
    \end{description}
    
    Dois sistema lineares são ditos \textbf{Sistemas Equivalentes} se e somente se admitem a mesma solução.
\end{teoria}

%------------------------------------------------------------------------------%
\begin{teoria}{Forma Matricial}
    
    Sempre podemos escrever um sistema linear em sua forma matricial 
    \[\left[
        \begin{array}{cccc}
            a_{11} & a_{12}  & \cdots & a_{1n} \\ 
            a_{21} & a_{22}  & \cdots & a_{2n} \\ 
            \vdots & \vdots  &        & \vdots \\
            a_{m1} & a_{m2}  & \cdots & a_{mn}
        \end{array} 
        \right]
        \left[
        \begin{array}{c} x_1 \\ x_2 \\ \vdots \\ x_n
        \end{array} 
        \right]
        =                    
        \left[
        \begin{array}{c} b_1 \\ b_2 \\ \vdots \\ b_m
        \end{array} 
        \right]
    \]    
    Abreviadamente podemos escrever $AX=B$ onde 
    \[A=\left[
        \begin{array}{cccc}
            a_{11} & a_{12}  & \cdots & a_{1n} \\ 
            a_{21} & a_{22}  & \cdots & a_{2n} \\ 
            \vdots & \vdots  &        & \vdots \\
            a_{m1} & a_{m2}  & \cdots & a_{mn}
        \end{array} 
        \right]
        \qquad X=
        \left[
        \begin{array}{c} x_1 \\ x_2 \\ \vdots \\ x_n
        \end{array} 
        \right]
        \qquad B=                    
        \left[
        \begin{array}{c} b_1 \\ b_2 \\ \vdots \\ b_m
        \end{array} 
        \right]
    \]
    
    Um sistema linear com o mesmo número de equações e incógnitas, $m=n$, será determinado
    se, e somente se, a sua matriz $A$ for inversível, isso é $\det A \neq 0$.

\end{teoria}

%------------------------------------------------------------------------------%
\begin{teoria}{Sistemas Lineares Homogêneos}
    
    Sistema lineares homogêneos são sistemas onde os termos independentes, $b_i$ são todos nulos.
    \[\left\{ 
    \begin{array}{ccccccccl}
    	a_{11}x_1 & \!+\! & a_{12}x_2 & \!+\! & \cdots & \!+\! & a_{1n}x_n & \!=\! & 0    \\
    	a_{21}x_1 & \!+\! & a_{22}x_2 & \!+\! & \cdots & \!+\! & a_{2n}x_n & \!=\! & 0    \\
    	 \vdots   &       &  \vdots   &       &        &       &  \vdots   &       & \vdots \\
    	a_{m1}x_1 & \!+\! & a_{m2}x_2 & \!+\! & \cdots & \!+\! & a_{mn}x_n & \!=\! & 0
    \end{array} 
    \right.
    \]    
    ou matricialmente 
    \[ AX = 0 \]
\end{teoria}

%------------------------------------------------------------------------------%
\begin{observacao}
    Não deixe de conferir suas respostas. Uma pequena distração, como uma troca de sinal,
    é suficiente para produzir falsos resultados.
\end{observacao}

%------------------------------------------------------------------------------%
\begin{teoria}{Regra de Cramer}
    
    A \textbf{Regra de Cramer} é um método para resolver um sistema linear, com $m=n$, determinado. A solução do sistema é calculada pela formula
    \[  x_i = \frac{D_i}{\det A} \qquad i=1,\dots,n \]
    onde $D_i$ é o determinante da matriz construída a partir da matriz $A$ substituindo a coluna $i$ pelos valores de $B$, por exemplo,
    \[D_2 = \left[
    \begin{array}{cccc}
    a_{11} & b_{1}  & \cdots & a_{1n} \\ 
    a_{21} & b_{2}  & \cdots & a_{2n} \\ 
    \vdots & \vdots &        & \vdots \\
    a_{n1} & b_{n}  & \cdots & a_{nn}
    \end{array} 
    \right]
     \]

    A Regra de Cramer tem grande importância teórica, mas existem métodos que demandam muito menos operações para chegar a solução de um sistema linear. 
    Principalmente quando $n$ é grande.
     
\end{teoria}

%------------------------------------------------------------------------------%
\section{Método de Gauss}

%------------------------------------------------------------------------------%
\begin{teoria}{Sistema Triangular Superior}
    
    Se a matriz de um sistema linear, com $m=n$, tiver todas as entradas abaixo da diagonal nulas, como nesse exemplo com $n=3$
    \[ A = \left[
        \begin{array}{ccc}
            a_{11} & a_{12} & a_{13} \\ 
            0      & a_{22} & a_{23} \\ 
            0      & 0      & a_{33}
        \end{array} 
        \right]
    \]    
    Dizermos que o sistema é \textbf{Triangular Superior}.
    \\
    
    Esse tipo de sistema pode ser facilmente resolvido se resolvermos primeiro a ultima equação encontrando $x_3$.
    Então substituímos esse valor na segunda equação e calculamos $x_2$. Por fim substituímos $x_3$ e $x_2$ na primeira
    equação e encontramos $x_1$. O mesmo método pode ser usado para resolver sistemas triangulares de qualquer dimensão.
    
\end{teoria}

%------------------------------------------------------------------------------%
\begin{teoria}{Método de Gauss}

    Como resolver um sistema triangular é relativamente fácil o \textbf{Método de Gauss} busca transformar
    o sistema que desejamos resolver em um sistema equivalente que seja triangular superior.
    \\

    Para transformar um sistema qualquer em um sistema triangular superior equivalente podemos usar qualquer uma das seguintes operações elementares nas linhas da matriz do sistema
    \begin{enumerate}
        \item Trocar duas linhas entre si.
        \item Multiplicar todos os elementos de uma linha por uma constante não-nula.
        \item Substituir uma linha pela sua soma com um múltiplo de outra.
    \end{enumerate}
        
    O Método de Gauss usa esses passos em uma sequencia específica que garante a transformação de qualquer matriz não singular em triangular superior.
        
\end{teoria}

%------------------------------------------------------------------------------%
\begin{exemplos}[2]
       
    \begin{exemplo}
        Para ilustrar a aplicação do Método de Gauss apresentamos aqui um exemplo com uma matriz de ordem 3.
        \begin{enumerate}
            \item Construir a matriz aumentada do sistema $[ A | B ]$
                \[ \left[
                \begin{array}{ccc|c}
                   	a_{11} & a_{12} & a_{13} & b_{1} \\
                   	a_{21} & a_{22} & a_{23} & b_{2} \\
                   	a_{31} & a_{32} & a_{33} & b_{3}
                \end{array} 
                \right]
                \]
            \item Queremos agora transformar anular os elementos abaixo da diagonal da matriz $A$. 
                Primeiro dividimos todos os elementos da primeira linha por $a_{11}$, 
                isso é, $L_1^{(1)} = L_1^{(0)} / a_{11}$, obtermos assim
                \[ \left[
                \begin{array}{ccc|c}
                 	  1    & a_{12}^{(1)} & a_{13}^{(1)} & b_{1}^{(1)} \\
                   	a_{21} &    a_{22}    &    a_{23}    &    b_{2}    \\
                   	a_{31} &    a_{32}    &    a_{33}    &    b_{3}
                \end{array} 
                \right] \]
            \item Substituímos agora a segunda linha por ela mesma subtraída da primeira linha multiplicada 
                por $a_{21}$, isso é, $L_2^{(1)} = L_2^{(0)} - L_1^{(1)}a_{21}$
                \[ \left[
                \begin{array}{ccc|c}
                 	  1    & a_{12}^{(1)} & a_{13}^{(1)} & b_{1}^{(1)} \\
                 	  0    & a_{22}^{(1)} & a_{23}^{(1)} & b_{2}^{(1)} \\
                   	a_{31} &    a_{32}    &    a_{33}    &    b_{3}
                \end{array} 
                \right] \]
            \item Fazemos o procedimento equivalente para anular o termo $a_{31}$, 
                isso é, $L_3^{(1)} = L_3^{(0)} - L_1^{(1)}a_{31}$
                \[ \left[
                \begin{array}{ccc|c}
                   	1 & a_{12}^{(1)} & a_{13}^{(1)} & b_{1}^{(1)} \\
                   	0 & a_{22}^{(1)} & a_{23}^{(1)} & b_{2}^{(1)} \\
                   	0 & a_{32}^{(1)} & a_{33}^{(1)} & b_{3}^{(1)}
                \end{array} 
                \right] \]
            \item Repetimos o processo para a segunda coluna. Primeiro dividimos a segunda 
                linha por $a_{22}^{(1)}$, isso é, $L_2^{(2)} = L_2^{(1)} / a_{22}^{(1)}$
                \[ \left[
                \begin{array}{ccc|c}
                   	1 & a_{12}^{(1)} & a_{13}^{(1)} & b_{1}^{(1)} \\
                   	0 &      1       & a_{23}^{(2)} & b_{2}^{(2)} \\
                   	0 & a_{32}^{(1)} & a_{33}^{(1)} & b_{3}^{(1)}
                \end{array} 
                \right] \]
            \item Substituímos agora a terceira linha por ela mesma subtraída da segunda 
                linha multiplicada por $a_{32}$, isso é, $L_3^{(2)} = L_3^{(1)} - L_2^{(1)}a_{32}$
                \[ \left[
                \begin{array}{ccc|c}
                   	1 & a_{12}^{(1)} & a_{13}^{(1)} & b_{1}^{(1)} \\
                   	0 &      1       & a_{23}^{(2)} & b_{2}^{(2)} \\
                   	0 &      0       & a_{33}^{(2)} & b_{3}^{(2)}
                \end{array} 
                \right] \]
            \item Dividimos agora a terceira linha por $a_{33}^(2)$, isso é, $L_3^{(3)} = L_3^{(2)} / a_{33}^{(2)}$
                \[ \left[
                \begin{array}{ccc|c}
                   	1 & a_{12}^{(1)} & a_{13}^{(1)} & b_{1}^{(1)} \\
                   	0 &      1       & a_{23}^{(2)} & b_{2}^{(2)} \\
                   	0 &      0       &      1       & b_{3}^{(3)}
                \end{array} 
                \right] \]
        \end{enumerate}
        Obtemos assim um sistema equivalente ao original que pode ser resolvido com muito mais facilidade
    \end{exemplo}

    \begin{exemplo}
        Resolva o sistema linear usando o Método de Gauss
        \[\systeme{
            	x  + y  + z  = 10,
            	2x + y  + 4z = 20,
            	2x + 3y + 5z = 25
        }\]
                
        Construindo a matriz aumentada do sistema
        \[\left[
            \begin{array}{rrr|r}
                1 & 1 & 1 & 10 \\ 
                2 & 1 & 4 & 20 \\ 
                2 & 3 & 5 & 25
            \end{array} 
            \right]
        \]
        Como o elemento $a_{11}$ já é igual a 1 não precisamos dividir a primeira linha por ele. 
        Subtraímos a segunda linha menos a primeira vezes 2 e a terceira linha menos a primeira, também, vezes 2
        \[\left[
            \begin{array}{rrr|r}
                1 &  1 & 1 & 10 \\ 
                0 & -1 & 2 & 0 \\ 
                0 &  1 & 3 & 5
            \end{array} 
            \right]
        \]
        Dividimos a segunda linha por $a_{22}=-1$ e obtemos
        \[\left[
            \begin{array}{rrr|r}
                1 &  1 &  1 & 10 \\ 
                0 &  1 & -2 & 0  \\ 
                0 &  1 &  3 & 5
            \end{array} 
            \right]
        \]
        Subtraímos a terceira linha menos a segunda linha
        \[\left[
            \begin{array}{rrr|r}
                1 &  1 &  1 & 10 \\ 
                0 &  1 & -2 & 0  \\ 
                0 &  0 &  5 & 5
            \end{array} 
            \right]
        \]
        Dividimos a terceira linha por $a_{33}=5$
        \[\left[
            \begin{array}{rrr|r}
                1 &  1 &  1 & 10 \\ 
                0 &  1 & -2 & 0  \\ 
                0 &  0 &  1 & 1
            \end{array} 
            \right]
        \]
        Resolvemos a terceira equação
        \[ \makebox[0pt][l]{$1z = 1$}\phantom{x + y + z = 10} \qquad \Rightarrow \qquad z = 1  \]
        Substituímos esse valor na segunda equação e a resolvemos também
        \[ \makebox[0pt][l]{$1y - 2z = 0$}\phantom{x + y + z = 10} \qquad \Rightarrow \qquad y = 2z = 2 \]
        Agora resolvemos a primeira equação
        \[ x + y + z = 10 \qquad \Rightarrow \qquad x = 10 - y - z = 10 -2 -1 = 7 \]
        A solução do sistema é
        \[ x = 7 \qquad y = 2 \qquad z = 1 \]
    \end{exemplo}
                            
\end{exemplos}

%------------------------------------------------------------------------------%
%------------------------------------------------------------------------------%
\section{Exercícios}

\openAnswers{03-Matrizes/03-Matrizes-Exercicios-ans}

\vspace{\baselineskip}
%------------------------------------------------------------------------------%
% \chapter{Matrizes}
%------------------------------------------------------------------------------%

%------------------------------------------------------------------------------%
\begin{exercicios}{Matrizes}
    
    \begin{exercicio}
        Indicar explicitamente os elementos da matriz $A=(a_{ij})_{2\times 3}$ tal que $a_{ij} = i - j$
        \begin{answers}
            \(\qquad A = \left[
                \begin{array}{rrr}
                	0 & -1 & -2 \\
                	1 &  0 & -1
                \end{array}
                \right]
            \)
        \end{answers}
    \end{exercicio}
    
    \begin{exercicio}
        Ache a matriz $A$ do tipo $2 \times 3$ definida por $a_{ij} = i j$ onde $i$ indica a linha e $j$, a coluna.
        \begin{answers}
            \( \qquad A = \left[
            \begin{array}{rrr}
            	1 & 2 & 3 \\
            	2 & 4 & 6
            \end{array}
            \right]
            \)
        \end{answers}
    \end{exercicio}
    
    \begin{exercicio}
        Construir as seguintes matrizes
        \begin{alphalist}
            \item $A=(a_{ij})_{3\times 3}$ tal que 
                  \( a_{ij} =\begin{cases}
                                1, & \text{se } i=j \\
                                0, & \text{se } i\neq j
                            \end{cases}  \)
            \item $B=(b_{ij})_{3\times 3}$ tal que 
                  \( b_{ij} =\begin{cases}
                                1, & \text{se } i+j=4 \\
                                0, & \text{se } i+j\neq 4
                            \end{cases}  \)
        \end{alphalist}
        \begin{answers}
            \begin{mathlist}[2]
                \item A = 
                    \left[
                    \begin{array}{rrr}
                    	1 & 0 & 0 \\
                    	0 & 1 & 0 \\
                    	0 & 0 & 1
                    \end{array}
                    \right]
                \item A = 
                    \left[
                    \begin{array}{rrr}
                    	0 & 0 & 1 \\
                    	0 & 1 & 0 \\
                    	1 & 0 & 0
                    \end{array}
                    \right]
            \end{mathlist}
        \end{answers}
    \end{exercicio}
    
    \begin{exercicio}
        Construa a matriz $A = (a_{ij})_{2 \times 3}$ definida por
        \[a_{ij} = \begin{cases}
            3i + j,  & \text{se } i < j \\
            7,       & \text{se } i = j \\
            i^2 + j, & \text{se } i > j
        \end{cases}\]
    \end{exercicio}
    
    \begin{exercicio}
        Dada as matrizes
        \[A = \left[\begin{array}{cc}
            x + y &   1   \\
             -5   & x - y
        \end{array} \right]
        \qquad
        B = \left[\begin{array}{rr}
             3 &  1 \\
            -5 & -1
        \end{array} \right]\]
        calcule $x$ e $y$ de modo que $A=B$.
    \end{exercicio}
        
\end{exercicios} % Matrizes

%------------------------------------------------------------------------------%
\begin{exercicios}{Operações com Matrizes}
        
    \begin{exercicio}
        Datas as matrizes 
        \[ A = \left[\begin{array}{cc}
            5 & 6 \\
            4 & 2
        \end{array} \right]
        \qquad
        B = \left[\begin{array}{rr}
            0 & -1 \\
            5 & 4
        \end{array} \right] \]
        calcule $A+B$, $A-B$, $2A+3B$, $AB$, $BA$, $A^t$ e $A^t+B$.
        \begin{answers}
            \begin{mathlist}[2]
                \item A+B   =   \left[
                \begin{array}{rr}
                    5 &  5 \\
                    9 &  6 \\
                \end{array}
                \right]
                \item A-B   =   \left[
                \begin{array}{rr}
                    5 &  7 \\
                    -1 & -2 \\
                \end{array}
                \right]
                \item 2A+3B =   \left[
                \begin{array}{rr}
                    10 &  9 \\
                    23 & 16 \\
                \end{array}
                \right]
                \item AB    =   \left[
                \begin{array}{rr}
                    30 & 19 \\
                    10 &  4 \\
                \end{array}
                \right]
                \item BA    =   \left[
                \begin{array}{rr}
                    -4 & -2 \\
                    41 & 38 \\
                \end{array}
                \right]
                \item A^t   =   \left[
                \begin{array}{rr}
                    5 &  4 \\
                    6 &  2 \\
                \end{array}
                \right]
                \item A^t+B =   \left[
                \begin{array}{rr}
                    5 &  3 \\
                    11 &  6 \\
                \end{array}
                \right]
            \end{mathlist}
        \end{answers}
    \end{exercicio}
    
    \begin{exercicio}
        Datas as matrizes
        \[ A = \left[\begin{array}{rrr}
            1 & 5 & 7 \\
            3 & 9 & 11
        \end{array} \right]
        \]\[
        B = \left[\begin{array}{rrr}
            2 &  4 &  6 \\
            8 & 10 & 12
        \end{array} \right]
        \]\[
        C = \left[\begin{array}{rrr}
            0 & -1 & -5 \\
            1 &  4 &  7
        \end{array} \right] \]
        calcule, se possível, $A+B+C$, $A-B+C$, $A-B-C$, $-A+B-C$, $AB$, $BC$, $AB^t$, $CB^t$, $A^tB$ e $A^tB^t$.
        \begin{answers}
            \begin{mathlist}[2]
            \item A+B+C    =   \left[
            \begin{array}{rrr}
                3 &  8 &  8 \\
                12 & 23 & 30 \\
            \end{array}
            \right]
            \item A-B+C    =   \left[
            \begin{array}{rrr}
                -1 &  0 & -4 \\
                -4 &  3 &  6 \\
            \end{array}
            \right]
            \item A-B-C    =   \left[
            \begin{array}{rrr}
                -1 &  2 &  6 \\
                -6 & -5 & -8 \\
            \end{array}
            \right]
            \item -A+B-C   =   \left[
            \begin{array}{rrr}
                1 &  0 &  4 \\
                4 & -3 & -6 \\
            \end{array}
            \right]
            \item AB \) Não pode ser calculada
            \item BC \) Não pode ser calculada
            \item AB^t     =   \left[
            \begin{array}{rr}
                64 & 142 \\
                108 & 246 \\
            \end{array}
            \right]
            \item CB^t     =   \left[
            \begin{array}{rr}
                -34 & -70 \\
                60 & 132 \\
            \end{array}
            \right]
            \item A^tB     =   \left[
            \begin{array}{rrr}
                26 & 34 & 42 \\
                82 & 110 & 138 \\
                102 & 138 & 174 \\
            \end{array}
            \right]
            \item A^tB^t \) Não pode ser calculada
            \end{mathlist}
        \end{answers}
    \end{exercicio}

    \begin{exercicio}
        Calcule os seguintes produtos matriciais
        \begin{mathlist}
            \item \left[\begin{array}{rrr}
                  1 & 2 & 3 \\
                  4 & 5 & 6
                \end{array} \right] 
                \left[\begin{array}{r}
                    7 \\ 8 \\ 9
                \end{array}\right]
            \item \left[\begin{array}{rr}
                   0 & 1 \\
                   1 & 0
                \end{array} \right] 
                \left[\begin{array}{rr}
                   4 & 7 \\ 
                   2 & 3
                \end{array}\right]
            \item \left[\begin{array}{rrr}
                    1 & 5 & 2 \\
                   -1 & 4 & 7
                \end{array} \right] 
                \left[\begin{array}{rr}
                    1 & -1 \\ 
                    2 &  3 \\
                   -3 &  0
                \end{array}\right]
            \item \left[\begin{array}{r}
                    1 \\ 2 \\ 3 
                \end{array} \right] 
                \left[\begin{array}{rrrr}
                    3 & 1 & 1 & 2
                \end{array}\right]
            \item \left[\begin{array}{rrrr}
                   1 & -1 & 5 & 0 \\
                   2 &  3 & 7 & 1 
                \end{array} \right] 
                \left[\begin{array}{rr}
                   1 & 1 \\ 
                   2 & 1 \\
                   3 & 1 \\
                   1 & 1
                \end{array}\right]
            \item \left[\begin{array}{rrr}
                   0 & 1 & 1 \\
                   2 & 2 & 0 \\
                   0 & 3 & 4 
                \end{array} \right] 
                \left[\begin{array}{rrr}
                   1 & 4 & 7 \\ 
                   0 & 0 & 1 \\
                   1 & 2 & 0
                \end{array}\right]
        \end{mathlist}
        \begin{answers}
            \begin{mathlist}[2]
                \item   \left[
                \begin{array}{r}
                    50 \\
                    122 \\
                \end{array}
                \right]
                \item   \left[
                \begin{array}{rr}
                    2 &  3 \\
                    4 &  7 \\
                \end{array}
                \right]
                \item   \left[
                \begin{array}{rr}
                    5 & 14 \\
                    -14 & 13 \\
                \end{array}
                \right]
                \item   \left[
                \begin{array}{rrrr}
                    3 &  1 &  1 &  2 \\
                    6 &  2 &  2 &  4 \\
                    9 &  3 &  3 &  6 \\
                \end{array}
                \right]
                \item   \left[
                \begin{array}{rr}
                    14 &  5 \\
                    30 & 13 \\
                \end{array}
                \right]
                \item   \left[
                \begin{array}{rrr}
                    1 &  2 &  1 \\
                    2 &  8 & 16 \\
                    4 &  8 &  3 \\
                \end{array}
                \right]
            \end{mathlist}
        \end{answers}
    \end{exercicio}

    \begin{exercicio}
        Sendo \( A=\left[\begin{array}{rr}
                                    1 & 1 \\
                                    0 & 1
                                 \end{array}\right] \)
        calcule $A^2$, $A^3$, $A^4$ e $A^n$.
        \begin{answers}
            \begin{mathlist}[2]
                \item A^2 =   \left[
                \begin{array}{rr}
                    1 &  2 \\
                    0 &  1 \\
                \end{array}
                \right]
                \item A^3 =   \left[
                \begin{array}{rr}
                    1 &  3 \\
                    0 &  1 \\
                \end{array}
                \right]
                \item A^4 =   \left[
                \begin{array}{rr}
                    1 &  4 \\
                    0 &  1 \\
                \end{array}
                \right]
                \item A^n =   \left[
                \begin{array}{rr}
                    1 &  n \\
                    0 &  1 \\
                \end{array}
                \right]
            \end{mathlist}
        \end{answers}
    \end{exercicio}
        
    \begin{exercicio}
        Calcule $AB$, $BA$, $A^2$ e $B^2$ sabendo que
        \[ A=\left[\begin{array}{rr}
                     2 &  1 \\
                    -4 & -2
                    \end{array}\right]
           \qquad
           B=\left[\begin{array}{rr}
                    2 & 1 \\
                    1 & 0
                   \end{array}\right]
        \]
        \begin{answers}
            \begin{mathlist}[2]
                \item AB  =   \left[
                \begin{array}{rr}
                    5 &  2 \\
                    -10 & -4 \\
                \end{array}
                \right]
                \item BA  =   \left[
                \begin{array}{rr}
                    0 &  0 \\
                    2 &  1 \\
                \end{array}
                \right]
                \item A^2 =   \left[
                \begin{array}{rr}
                    0 &  0 \\
                    0 &  0 \\
                \end{array}
                \right]
                \item B^2 =   \left[
                \begin{array}{rr}
                    5 &  2 \\
                    2 &  1 \\
                \end{array}
                \right]
            \end{mathlist}
        \end{answers}
    \end{exercicio}
        
    \begin{exercicio}
        Calcular o produto $ABC$ sendo dados
        \[   A=\left[\begin{array}{rr}
                 1 & 2 \\
                 5 & 1
                \end{array}\right]
            \hfill
            B=\left[\begin{array}{rrr}
                1 & 1 & 1\\
                3 & 2 & 1
                \end{array}\right]
            \hfill
            C=\left[\begin{array}{rr}
                3 &  1 \\
                1 &  0 \\
                2 & -1
                \end{array}\right]
        \]
        \begin{answers}
            \[ ABC =   \left[
                \begin{array}{rr}
                	32 & 4 \\
                	43 & 2
                \end{array}
                \right]
            \]
        \end{answers}
    \end{exercicio}
    
   \begin{exercicio}
       Considerando as matrizes
       \[ A = \left[
       \begin{array}{rrr}
       1 & 2 & -3 \\
       3 & 4 &  0
       \end{array} 
       \right]
       \qquad
       B = \left[
       \begin{array}{rrr}
       -2 & 1 &  5 \\
       0 & 3 & -4
       \end{array} 
       \right]
       \]
       Calcule
       \begin{mathlist}[2]
           \item C = A + B
           \item D = A - B
           \item E = 2A + 3B
           \item F = A^t + B^t
           \item G = A B^t
           \item H = A^t B
        \end{mathlist}
        \begin{answers}
            \begin{mathlist}[2]
                \item C =   \left[
                \begin{array}{rrr}
                    -1 &  3 &  2 \\
                    3 &  7 & -4 \\
                \end{array}
                \right]
                \item D =   \left[
                \begin{array}{rrr}
                    3 &  1 & -8 \\
                    3 &  1 &  4 \\
                \end{array}
                \right]
                \item E =   \left[
                \begin{array}{rrr}
                    -4 &  7 &  9 \\
                    6 & 17 & -12 \\
                \end{array}
                \right]
                \item F =   \left[
                \begin{array}{rr}
                    -1 &  3 \\
                    3 &  7 \\
                    2 & -4 \\
                \end{array}
                \right]
                \item G =   \left[
                \begin{array}{rr}
                    -15 & 18 \\
                    -2 & 12 \\
                \end{array}
                \right]
                \item H =   \left[
                \begin{array}{rrr}
                    -2 & 10 & -7 \\
                    -4 & 14 & -6 \\
                    6 & -3 & -15 \\
                \end{array}
                \right]
            \end{mathlist}
        \end{answers}
    \end{exercicio}
    
    \begin{exercicio}
        Considerando as matrizes
        \[ A = \left[ 
            \begin{array}{rrr}
               	1 & 2 & -3 \\
               	3 & 4 &  0
            \end{array} 
            \right]
            \qquad
            B = \left[
            \begin{array}{rrr}
               	-2 &  1 & 0 \\
              	 0 &  3 & 0 \\
              	 5 & -4 & 0
            \end{array} 
            \right]
        \]
        Calcule
        \begin{mathlist}[2]
            \item C = AB
            \item D = A(3B)
            \item E = B^2
            \item F = B^3
            \item G = B B^t
            \item H = A^t A
            \item J = A A^t
        \end{mathlist}
        \begin{answers}
            \begin{mathlist}[2]
                \item C =   \left[
                \begin{array}{rrr}
                    -17 & 19 &  0 \\
                    -6 & 15 &  0 \\
                \end{array}
                \right]
                \item D =   \left[
                \begin{array}{rrr}
                    -51 & 57 &  0 \\
                    -18 & 45 &  0 \\
                \end{array}
                \right]
                \item E =   \left[
                \begin{array}{rrr}
                    4 &  1 &  0 \\
                    0 &  9 &  0 \\
                    -10 & -7 &  0 \\
                \end{array}
                \right]
                \item F =   \left[
                \begin{array}{rrr}
                    -8 &  7 &  0 \\
                    0 & 27 &  0 \\
                    20 & -31 &  0 \\
                \end{array}
                \right]
                \item G =   \left[
                \begin{array}{rrr}
                    5 &  3 & -14 \\
                    3 &  9 & -12 \\
                    -14 & -12 & 41 \\
                \end{array}
                \right]
                \item H =   \left[
                \begin{array}{rrr}
                    10 & 14 & -3 \\
                    14 & 20 & -6 \\
                    -3 & -6 &  9 \\
                \end{array}
                \right]
                \item J =   \left[
                \begin{array}{rr}
                    14 & 11 \\
                    11 & 25 \\
                \end{array}
                \right]
            \end{mathlist}
        \end{answers}
    \end{exercicio}
    
    \begin{exercicio}
        Considere as seguintes matrizes
        \[
        A = \left[\begin{array}{cc}
           	2 & 0 \\
           	6 & 7
        \end{array}\right]
        \qquad
        B = \left[\begin{array}{rr}
           	0 &  4 \\
           	2 & -8
        \end{array}\right]
        \]\[
        C = \left[\begin{array}{rrr}
           	-6 &  9 & -7 \\
          	 7 & -3 & -2
        \end{array}\right]
        \]\[ 
        D = \left[\begin{array}{rrr}
           	-6 & 4 & 0 \\
          	 1 & 1 & 4 \\
           	-6 & 0 & 6
        \end{array}\right] 
        \qquad
        E = \left[\begin{array}{rrr}
          	 6 & 9 & -9 \\
           	-1 & 0 & -4 \\
           	-6 & 0 & -1
        \end{array}\right] 
        \]
        Se for possível calcule
        \begin{mathlist}[2]
            \item AB-BA
            \item 2C-D
            \item (2D^t-3E^t)^t
            \item D^2-DE
        \end{mathlist}
        \begin{answers}
            \begin{mathlist}[2]
                \item AB-BA         =   \left[
                \begin{array}{rr}
                    -24 & -20 \\
                    58 & 24 \\
                \end{array}
                \right]
                \item 2C-D \) Não pode ser calculada
                \item (2D^t-3E^t)^t =   \left[
                \begin{array}{rrr}
                    -30 & -19 & 27 \\
                    5 &  2 & 20 \\
                    6 &  0 & 15 \\
                \end{array}
                \right]
                \item D^2-DE        =   \left[
                \begin{array}{rrr}
                    80 & 34 & -22 \\
                    -10 & -4 & 45 \\
                    72 & 30 & -12 \\
                \end{array}
                \right]
            \end{mathlist}
        \end{answers}
    \end{exercicio}    

            
    \begin{exercicio}
        Sendo 
        \[ A=\left[\begin{array}{rr}
        	1 & 9 \\
        	2 & 1
        \end{array}\right]
        \qquad
        B=\left[\begin{array}{rr}
        	0 & -8 \\
        	3 & -1
        \end{array}\right]\]
        calcule
        \begin{mathlist}[2]
            \item (A+B)^2
            \item (A+B)(A-B)
            \item A^2 -2I_2A +I_2^2
            \item A^3 - I_2^3
        \end{mathlist}
        \begin{answers}
            \begin{mathlist}
                \item (A+B)^2           =   \left[
                \begin{array}{rr}
                    6 &  1 \\
                    5 &  5 \\
                \end{array}
                \right]
                \item (A+B)(A-B)        =   \left[
                \begin{array}{rr}
                    0 & 19 \\
                    5 & 85 \\
                \end{array}
                \right]
                \item A^2 -2I_2A +I_2^2 =   \left[
                \begin{array}{rr}
                    18 &  0 \\
                    0 & 18 \\
                \end{array}
                \right]
                \item A^3 - I_2^3       =   \left[
                \begin{array}{rr}
                    54 & 189 \\
                    42 & 54 \\
                \end{array}
                \right]
            \end{mathlist}
        \end{answers}
    \end{exercicio}
            
    \begin{exercicio}
        Sendo \( A=\left[\begin{array}{rr}
                                    1 & -1 \\
                                    0 &  2
                                 \end{array}\right] \),
        qual das matrizes comuta com $A$
        \begin{mathlist}[2]
            \item B = \left[\begin{array}{r}
                    2 \\ 3
                \end{array}\right]
            \item C = \left[\begin{array}{rrr}
                   1 & 3 & 2 \\
                   4 & 5 & 1
                \end{array}\right]
            \item D = \left[\begin{array}{rr}
                   0 & 0 \\
                   1 & 0
                \end{array}\right]
            \item E = \left[\begin{array}{rr}
                   5 & 2 \\
                   0 & 3
                \end{array}\right]
        \end{mathlist}
    \end{exercicio}
        
    \begin{exercicio}
        Calcule $x$, $y$, $z$ e $t$ sabendo que
        \[\left[\begin{array}{cc}
        x & 1 \\
        1 & 2
        \end{array}\right] +
        \left[\begin{array}{rr}
        2 &  y \\
        0 & -1
        \end{array}\right]  = 
        \left[\begin{array}{cc}
        3 & 2 \\
        z & t
        \end{array}\right]\]
    \end{exercicio}
    
    \begin{exercicio}
        Obtenha $X$ sabendo que 
        \[ X + \left[\begin{array}{r} 1 \\  4 \\  7 \end{array} \right]
        = \left[\begin{array}{r} 5 \\  7 \\  2 \end{array} \right]
        + \left[\begin{array}{r} 1 \\ -1 \\ -2 \end{array} \right]
        \]
    \end{exercicio}        
        
    \begin{exercicio}
        Determine $x$ e $y$ de modo que mas matrizes $A$ e $B$ comutem
        \[ A=\left[\begin{array}{rr}
                1 & 2 \\
                1 & 0
            \end{array}\right]
            \qquad
            B=\left[\begin{array}{rr}
                0 & 1 \\
                x & y
            \end{array}\right]
        \]
    \end{exercicio}
            
    \begin{exercicio}
        Obter todas as matrizes $B$ que comutem com 
          \[ A=\left[\begin{array}{rr}
                        1 & -1 \\
                        3 &  0
                    \end{array}\right]\]
    \end{exercicio}

    \begin{exercicio}
        Dada as matrizes
        \[A = \left[\begin{array}{r}
            -1 \\
             2
        \end{array} \right]
        \qquad
        B = \left[\begin{array}{c}
            3 \\
            0
        \end{array} \right]\] 
        \[C = \left[\begin{array}{r}
            -2 \\
             4
        \end{array} \right] 
        \qquad
        D = \left[\begin{array}{r}
             1 \\
            -1
        \end{array} \right]\]
        determine $X$ e $Y$, tais que
        \begin{mathlist}[1]
            \item X = A + 2B - C + D
            \item \frac{Y + A}{2} = \frac{C}{2} + 2B - 3D
        \end{mathlist}
    \end{exercicio}

    \begin{exercicio}
        Sendo
        \[A = \left[\begin{array}{cc}
            3 & 1 \\
            5 & 2
        \end{array} \right]
        \qquad
        B = \left[\begin{array}{cc}
            2 & 1 \\
            1 & 1
        \end{array} \right]\]
        \[I = \left[\begin{array}{cc}
            1 & 0 \\
            0 & 1
        \end{array} \right]    
        \qquad
        O = \left[\begin{array}{cc}
            0 & 0 \\
            0 & 0
        \end{array} \right]\]
        resolva o sistema, onde $X$ e $Y$ são matrizes quadradas de ordem 2
        \[\begin{cases}
            AX + Y = O \\
            BX + Y = I
        \end{cases}\]
    \end{exercicio}
    
    \begin{exercicio}
        Calcule $x$, data a igualdade
        \[\left[\begin{array}{cr}
            x & 1  \\
            3 & -1
        \end{array} \right]
        \left[\begin{array}{cc}
            2 & 0 \\
            1 & 1
        \end{array} \right]    =
        \left[\begin{array}{cr}
            0 & 1  \\
            5 & -1
        \end{array} \right]\]
    \end{exercicio}

    \begin{exercicio}
        Determine os valores de $x$, $y$ e $z$ na igualdade abaixo, envolvendo matrizes reais $2 \times 2$
        \[\left[\begin{array}{cc}
        0 & 0 \\
        x & 0
        \end{array} \right]    \,
        \left[\begin{array}{cc}
        0 & x \\
        0 & 0
        \end{array} \right] = 
        \left[\begin{array}{cc}
        x - y & 0 \\
        x   & z
        \end{array} \right] +
        \left[\begin{array}{cc}
        z - 4 & 0 \\
        y - z & 0
        \end{array} \right]\]
    \end{exercicio}
    
    \begin{exercicio}
        Calcule $x$ de modo que as matrizes
        \[A = \left[\begin{array}{rr}
        2 & -3 \\
        1 & 4
        \end{array} \right]    
        \qquad
        B = \left[\begin{array}{cc}
        4 & x \\
        0 & 4
        \end{array} \right]\]
        sejam comutativas.
    \end{exercicio}
    
    \begin{exercicio}
        Dada a matriz
        \[ A = \left[\begin{array}{cc}
        \sfrac{3}{4} & \sfrac{1}{2} \\[2mm]
        \sfrac{1}{4} & \sfrac{1}{2}
        \end{array} \right] \]
        determine uma matriz real da forma
        \[T = \left[\begin{array}{c}
        x   \\
        1 - x
        \end{array} \right]\]
        para a qual $AT = T$
    \end{exercicio}

    \begin{exercicio}
       Seja, para cada $x$ real, a matriz quadrada de ordem~2
       \[M (x) = \left[\begin{array}{rr}
            x & -x \\
           2x &  x
       \end{array} \right]\]
       Calcule o produto de $M(x)$ por $M(-x)$.
    \end{exercicio}
       
   \begin{exercicio}
       Se conhecemos apenas as matrizes $AB$ e $AC$ como podemos calcular as matrizes
       \begin{mathlist}[2]
           \item A(B+C)
           \item B^tA^t
           \item C^tA^t
           \item (ABA)C
       \end{mathlist}
    \end{exercicio}

    \begin{exercicio}
        Mostre que as matrizes da forma
        \[A = \left[ \begin{array}{ll}
            1 & y^{-1} \\
            y & 1
        \end{array} \right]\]
        em que $y$ é uma número real não nulo, verificam a equação \[X^2 = 2X\]
    \end{exercicio}
    
    \begin{exercicio}
        Verifique se as identidades são verdadeiras
        \begin{mathlist}
            \item (A+B)^2 = A^2 + 2AB + B^2
            \item (AB)C = C(AB)
            \item (ABC)^t = C^tB^tA^t
            \item (A+I)(A-I) = A^2 - I 
        \end{mathlist}
    \end{exercicio}

\end{exercicios} % Operações com Matrizes

%------------------------------------------------------------------------------%
\begin{exercicios}{Matriz Transposta}

    \begin{exercicio}
        Dada a matriz
        \[A = \left[\begin{array}{rrr}
        	 1 &  4 & 0 \\
        	-2 & -1 & 1 \\
        	 3 &  5 & 2
        \end{array} \right]    \]
        calcule
        \begin{mathlist}[2]
            \item X = A + A^t
            \item Y = A^t - A
        \end{mathlist}
        \begin{answers}
            \begin{mathlist}[2]
                \item X =   \left[
                \begin{array}{rrr}
                    2 &  2 &  3 \\
                    2 & -2 &  6 \\
                    3 &  6 &  4 \\
                \end{array}
                \right]
                \item Y =   \left[
                \begin{array}{rrr}
                	 0 & -6 & 3 \\
                	 6 &  0 & 4 \\
                	-3 & -4 & 0
                \end{array}
                \right]
            \end{mathlist}
        \end{answers}
    \end{exercicio}

    \begin{exercicio}
        Considere as matrizes
        \[A = \left[\begin{array}{rrr}
             a & b & 1 \\
            -1 & 1 & a
        \end{array} \right]
        \qquad
        B = \left[\begin{array}{rrr}
            1 & -1 & 0 \\
            0 &  1 & 0
        \end{array} \right]\]
        Determine $a$ e $b$, sabendo que
        \[A B^t = \left[\begin{array}{rr}
             3 & 4 \\
            -2 & 1
        \end{array} \right]\]
    \end{exercicio}

    \begin{exercicio}
        Sabendo que 
        \[ A = \left[ \begin{array}{r} x \\  4 \\ -2 \end{array} \right] \qquad
           B = \left[ \begin{array}{r} 2 \\ -3 \\  5 \end{array} \right] \]
        e que $A^tB=4$ determine o valor de $x$.
    \end{exercicio}

    \begin{exercicio}
        Encontre um valor de $x$ tal que $AB^t=0$, sendo que 
        \[ A = \left[ \begin{array}{ccc} x & 4 & -2 \end{array} \right]  \qquad
           B = \left[ \begin{array}{ccc} 2 & -3 & 5 \end{array} \right]
        \]
    \end{exercicio}
     
\end{exercicios} % Matriz Transposta
            
%------------------------------------------------------------------------------%
\begin{exercicios}{Matriz Inversa}
        
    \begin{exercicio}
        Mostre que $A^{-1}$ é a inversa de $A$
        \begin{mathlist}
            \item A = \left[\begin{array}{cc}
                     1 & 3 \\
                     2 & 7
                 \end{array} \right] 
                 \quad
                 A^{-1} = \left[\begin{array}{cc}
                     7  & -3 \\
                     -2 & 1
                 \end{array} \right]
            \item A = \left[\begin{array}{rrr}
                                        1 & 2 & 7 \\
                                        0 & 3 & 1 \\
                                        0 & 5 & 2
                                    \end{array} \right] \)
                  \( A^{-1} = \left[ \begin{array}{rrr}
                                            1 & 31 & -19 \\
                                            0 &  2 &  -1 \\
                                            0 & -5 &   3
                                        \end{array} \right] \)
        \end{mathlist}
    \end{exercicio}

    \begin{exercicio}
        Resolva a equação  matricial $AX = B$, dadas as matrizes
        \[A = \left[\begin{array}{rr}
             2 & -1 \\
            -3 &  2
        \end{array} \right]    
        \qquad
        B = \left[\begin{array}{rr}
             3 &  2 \\
            -6 & -5
        \end{array} \right]    \]
    \end{exercicio}

    \begin{exercicio}
        Determine $a$ real de modo que a matriz 
        \[A = \left[\begin{array}{cc}
            1 & 2 \\
            0 & a
        \end{array} \right]\]
        seja igual à sua inversa.
    \end{exercicio}
    
\end{exercicios} % Matriz Inversa

%------------------------------------------------------------------------------%

%------------------------------------------------------------------------------%
% \chapter{Matrizes}
%------------------------------------------------------------------------------%

%------------------------------------------------------------------------------%
\begin{exercicios}{Sistemas Lineares}
    
    \begin{exercicio}
        Diga qual das equações abaixo são lineares
        \begin{mathlist}
            \item x_1-x_2+x_3-2x_4 =3
            \item x_1 + mx_2+x_3^2=n \) \\[1mm] onde $m$ e $n$ são constantes dadas
            \item x-2y+3z=4
            \item a_1x_1 + a_2x_2^2+a_3x_3^3=b \) \\[1mm] onde $a_i$ e $b_i$ são constantes dadas
            \item 2x_1 + \log(x_2) +x_3 =\log(2)
            \item -x_1 +x_2 -3x_3 -x_4 - x_5 = 0
            \item \sqrt{3}x_1 +\sqrt{2}x_2 +x_3 = 5
            \item x_1 + 3x_2 - 4x_3 + 5x_4 = 10 - 2x_5
            \item e^2x_1 + \log(4)x_2 = \sqrt{3}
            \item x_1 + x_1x_2 = x_3
        \end{mathlist}
    \end{exercicio}
    
    \begin{exercicio}
        Verifique se $(2,0,-3)$ é solução de \[ 2x_1+5x_2+2x_3=-2 \]
    \end{exercicio}
    
    \begin{exercicio}
        Verifique se $(1,1,-1,-1)$ é solução de \[ 5x_1 -10x_2 -x_3 +2x_4=0 \]
    \end{exercicio}
        
    \begin{exercicio}
        Encontre uma solução para a equação linear 
        \[ 2x_1 -x_2 -x_3 = 0 \]
        diferente de $(0,0,0)$
    \end{exercicio}
        
    \begin{exercicio}
        Escreva os seguintes sistemas na forma matricial, 
        considerando que $a$, $b$, $c$, $d$, $e$, $f$, $m$ e $n$ são constantes
        \begin{mathlist}
            \item \systeme{
                x-y+z    =2,
            	-x+2y+2z =5,
            	5x-y+5z  =1}
            \item \systeme{
            	3x-5y+4z-t = 8,
            	2x+y-2z    = -3,
            	-x-2y+z-3t = 1,
            	-5x -y+6t  = 4}
            \item \systeme{
            	ax+by+cz   = d,
            	-mz + ny   = e,
            	abx-b^2+mz = f}
            \item \systeme{
            	\sqrt{2}x-3y+2z  = 7,
            	+7y -z           = 0,
            	4x+\sqrt{3}y +2z = 5}
            \item \systeme{
            	ax-by+2z  = 1,
            	a^2-by +z = 3}
            \item \systeme{
            	x+y-z   = 3-t,
            	-x-y-2z =1-3t,
            	5x+3z   =7+t}
        \end{mathlist}
    \end{exercicio}
    
    \begin{exercicio}
        Verifique se $(0,-3,-4)$ é solução do sistema
        \[\systeme{
        	x+y-z      =1,
        	2x-y+z     = -1,
        	x + 2y + z = 2}
        \]
    \end{exercicio}

    \begin{exercicio}
        Verifique se $(1,0,-2,1)$ é solução do sistema
        \[\systeme{
        	5x+3y-2z-4t  =5,
        	2x-4y +3z -5t =-9,
        	-x+2y-5z+3t   =12}
        \]
    \end{exercicio}
        
    \begin{exercicio}
        Encontre a solução do sistema $(A+4I_3)X=0$ sabendo que
        \[ A = \left[ 
        \begin{array}{rrr}
            1 & 0 & 5 \\
            1 & 1 & 1 \\
            0 & 1 & 4 \\
        \end{array}
        \right]\]
    \end{exercicio}
    
    \begin{exercicio}
        Encontre os valores de $a$ para os quais a matriz
        \[A = \left[ 
        \begin{array}{rrr}
            1 & 1 & 0 \\
            1 & 0 & 0 \\
            1 & 2 & a \\
        \end{array}
        \right]\]
        possui inversa.
    \end{exercicio}
    
    \begin{exercicio}
        Sejam
        \[A = \left[ 
            \begin{array}{rrr}
                1 & 0 & 5 \\
                1 & 1 & 1 \\
                0 & 1 & 4 \\
            \end{array}
            \right] 
            \qquad  
        B = \left[ \begin{array}{c} 5 \\ 5 \\ 10 \\ \end{array} \right] 
        \]
        encontre a solução do sistema $(A+4I_3)X=B$.
    \end{exercicio}
    
    \begin{exercicio}
        Resolva o sistema
        \[\begin{cases}
        	\dfrac{2}{u} + \dfrac{3}{v} =  8 \\[4mm]
        	\dfrac{1}{u} - \dfrac{1}{v} = -1
        \end{cases}\]
    \end{exercicio}
    
    \begin{exercicio}
        Os valores de $x$, $y$ e $z$, solução do sistema
        \[\systeme{
        	x + 2y + 3z  = 14,
        	4x + 5y + 6z = 32,
        	7x + 8y + 9z = a}
        \]
        formam, nesta ordem, uma P.A. de razão $1$. Qual é o valor de $a$?	
    \end{exercicio}
    
    \begin{exercicio}
        Resolva o sistema utilizando a regra de Cramer
        \[\systeme{
            4x - 3y = 11,
            2x + 5y = -1}
        \]
    \end{exercicio}
    
    \begin{exercicio}
        Resolva o sistema a seguir utilizando a regra de Cramer
        \[\systeme{
        	2x + y - z  = 5,
        	3x - 2y + z = -2,
        	x + z       = 0}
        \]
    \end{exercicio}
    
    \begin{exercicio}
        Discuta o sistema, em função dos valores do parâmetro real $a$.
        \[\systeme{
            a^2x + y = 1,
            x + y    = a}
        \]
    \end{exercicio}
    
    \begin{exercicio}
        Discuta, por Cramer, o sistema
        \[\systeme{
            7x + y - 3z = 10,
            x + y + z   = 6,
            4x + y + Pz = Q}
        \]
    \end{exercicio}
    
    \begin{exercicio}
        Encontre o valor de $a$ para que o sistema seja possível. 
        \[\systeme{
            2x - y + 3z  = a,
            x + 2y - z   = 3,
            7x + 4y + 3z = 13}
        \]
        Para o valor encontrado de $a$, ache a solução geral do sistema, isto é, 
        ache expressões que representem todas as soluções do sistema. 
        Explicite duas dessas soluções.		
    \end{exercicio}
    
    \begin{exercicio}
        Calcule o valor de $m$ para que o sistema admita soluções diferentes da trivial
        \[\systeme{
            x + y - z   = 0,
            2x - y + z  = 0,
            3x + my - z = 0}
        \]
    \end{exercicio}
    
    \begin{exercicio}
        Sejam $a$, $b$ e $c$ números tais que
        \[\systeme{
            a + 2b = 5,
            b + 2c = 8,
            2a + c = 5}
        \]
        Determine o valor de $a + b + c$.		
    \end{exercicio}
    
    \begin{exercicio}
        Calcule $a$ e $b$ de modo que o sistema a seguir seja indeterminado
        \[\systeme{
            x + 2y + 2z  = a,
            3x + 6y - 4z = 4,
            2x + by - 6z = 1}
        \]
    \end{exercicio}
    
    \begin{exercicio}
        Resolva o sistema
        \[\systeme{
            2y + x = b,
            2z - y = b,
            az + x = b}
        \]
        para os casos
        \begin{mathlist}[2]
            \item a = 0 \text{ e } b = 1
            \item a = 4 \text{ e } b = 0
        \end{mathlist}		
    \end{exercicio}
    
    \begin{exercicio}
        Ache os valores de $k \in \R$ de maneira que o sistema linear seja impossível.
        \[\systeme{
             x + y = 1,
             y + z = 2,
            kz + x = 0}
        \]
    \end{exercicio}
    
\end{exercicios} % Sistemas Lineares

%------------------------------------------------------------------------------%
\begin{exercicios}{Método de Gauss}

    \begin{exercicio}
        Resolva o sistema pelo método de Gauss
        \[ \systeme{  x + 3y = 7,
                     2x +  y = 4} 
        \]
        \begin{answers}
            \(\qquad x=1 \qquad y=2 \)
        \end{answers}
    \end{exercicio}
    
    \begin{exercicio}
        Resolva o sistema pelo método de Gauss
        \[\systeme{ 2x + y -2z = 10,
            3x +2y +2z =  1,
            5x +4y +3z =  4}
        \]
        \begin{answers}
            \(\qquad x=1 \qquad y=2 \qquad z=-3 \)
        \end{answers}
    \end{exercicio}
    
    \begin{exercicio}
        Resolva o seguinte sistema de equações lineares
        \[\systeme{2x +  y +  z +  w = 1,
            x + 2y +  z +  w = 2,
            x +  y + 2z +  w = 3,
            x +  y +  z + 2w = 4}
        \]
        \begin{answers}
            \(\qquad w=-1 \qquad x=0 \qquad y=1 \qquad z=2 \)
        \end{answers}
    \end{exercicio}
    
    \begin{exercicio}
        Resolva o sistema pelo método de Gauss
        \[\systeme{ x  -y -2z = 1,
            -x  +y + z = 2,
            x -2y + z =-2}
        \]
        \begin{answers}
            \(\qquad x=-11 \qquad y=-6 \qquad z=-3 \)
        \end{answers}
    \end{exercicio}
    
    \begin{exercicio}
        Resolva o sistema pelo método de Gauss
        \[\systeme{  x -2y +3z =  1,
            3x +6y -3z = -2,
            6x +6y +3z =  5}
        \]
        \begin{answers}
            \(\qquad x=2 \qquad y=\sfrac{7}{4} \qquad z=\sfrac{13}{6} \)
        \end{answers}
    \end{exercicio}
    
    \begin{exercicio}
        Use método de Gauss para resolver o sistema linear $(A+2I)X=B$, onde
        \[A = \left[ \begin{array}{rrr}
        1 & 0 & 5 \\
        1 & 1 & 1 \\
        0 & 1 & 4
        \end{array} \right] 
        \qquad
        B = \left[ \begin{array}{r} 3 \\ 4 \\ 1 \\ \end{array}
        \right] 
        \]
        \begin{answers}
            \(
            \qquad A+2I = \left[
            \begin{array}{rrr}
            3 & 0 & 5 \\
            1 & 3 & 1 \\
            0 & 1 & 6
            \end{array}
            \right]
            \qquad X = \left[
            \begin{array}{r}
            1 \\
            1 \\
            0
            \end{array}
            \right]
            \)
        \end{answers}
    \end{exercicio}
            
    \begin{exercicio}
        Resolva se possível o sistema pelo método de Gauss
        \[\systeme{ x - 2y +  z = 1,
                   2x +  y -  z = 2,
                    x + 3y - 2z = 1}
        \]
        \begin{answers}
            \(\qquad\) Sistema possui infinitas soluções.
        \end{answers}
    \end{exercicio}
    
    \begin{exercicio}
        Resolva se possível o sistema pelo método de Gauss
        \[\systeme{ x + 3y -  z = 6,
                   2x + 4y + 2z = 9,
                    x +  y + 3z = 1}
        \]
        \begin{answers}
            \(\qquad\) Sistema não possui solução.
        \end{answers}
    \end{exercicio}
    
    \begin{exercicio}
        Resolva se possível o sistema pelo método de Gauss
        \[\systeme{ 3x + 4y -  z = 1,
                    4x + 5y + 2z = 12,
                     x - 2y + 3z = 8}
        \]
        \begin{answers}
            \(\qquad x=0,1 \qquad y=1 \qquad z=3,3 \)
        \end{answers}
    \end{exercicio}
    
    \begin{exercicio}
        Resolva se possível o sistema pelo método de Gauss
        \[\systeme{  x + 2y -3z = 1,
            3x -  y +2z = 7,
            5x + 3y -4z = 2}
        \]
        \begin{answers}
            \(\qquad x=-7 \qquad y=79 \qquad z=50 \)
        \end{answers}
    \end{exercicio}
    
    \begin{exercicio}
        Resolva se possível o sistema pelo método de Gauss
        \[\systeme{  x +2y -3z =  6,
            2x - y +4z =  2,
            4x +3y -2z = 14}
        \]
        \begin{answers}
            \(\qquad x=\sfrac{7}{3} \qquad y=\sfrac{4}{3} \qquad z=-\sfrac{1}{3} \)
        \end{answers}
    \end{exercicio}
    
    \begin{exercicio}
        Resolva se possível o sistema pelo método de Gauss
        \[\systeme{ 2x +3y +z = 0,
            x -2y -z = 1,
            x +4y +z = 2}
        \]
        \begin{answers}
            \(\qquad x=-\sfrac{1}{4} \qquad y=\sfrac{7}{4} \qquad z=-\sfrac{19}{4} \)
        \end{answers}
    \end{exercicio}
    
    \begin{exercicio}
        Resolva se possível o sistema pelo método de Gauss
        \[\systeme{-x + y - 2z      = 1,
            2x - y +      3t = 2,
            x -2y +  z - 2t = 0}
        \]
        \begin{answers}
            \(\qquad \) Sistema possui infinitas soluções
        \end{answers}
    \end{exercicio}
    
    \begin{exercicio}
        Resolva se possível o sistema pelo método de Gauss
        \[\systeme{ x +3y  + 2z = 2,
            3x +5y +4z = 4,
            5x +3y +4z = -10}
        \]
        \begin{answers}
            \(\qquad \) Sistema não possui solução
        \end{answers}
    \end{exercicio}
    
    \begin{exercicio}
        Resolva se possível o sistema pelo método de Gauss
        \[\systeme{ x + y - z + t = 1,
            3x - y -2z + t = 2,
            -x -2y +3z +2t = -1}
        \]
        \begin{answers}
            \(\qquad \) Sistema possui infinitas soluções
        \end{answers}
    \end{exercicio}
    
    \begin{exercicio}
        Resolva se possível o sistema pelo método de Gauss
        \[\systeme{ x + y + z + t = 1,
            x - y + z + t = -1,
            y - z +2t = 2,
            2x + z - t = -1 }
        \]
        \begin{answers}
            \(\qquad x=-\sfrac{1}{5} \qquad y=1 \qquad z=-\sfrac{1}{5} \qquad t=\sfrac{2}{5} \)
        \end{answers}
    \end{exercicio}
    
    \begin{exercicio}
        Resolva se possível o sistema pelo método de Gauss
        \[\systeme{ x -2y -3z = 5,
            -2x +5y +2z = 3,
            -x +3y - z = 2}
        \]
        \begin{answers}
            \(\qquad \) Sistema não possui solução
        \end{answers}
    \end{exercicio}
    
\end{exercicios} % Método de Gauss

%------------------------------------------------------------------------------%

%------------------------------------------------------------------------------%
% \chapter{Matrizes}
%------------------------------------------------------------------------------%

%------------------------------------------------------------------------------%
\begin{exercicios}{Determinante de Uma Matriz e Suas Propriedades}
    
    
% 1
\begin{exercicio}
    Calcule os determinantes
    \begin{mathlist}[2]
        \item \left|\begin{array}{rr}
            -3 & -1 \\
            2 &  4
        \end{array} \right|
        \item \left|\begin{array}{rr}
            13 & 7 \\
            11 & 5
        \end{array} \right|
        \item \left|\begin{array}{rr}
            3 & 1 \\
            5 & 2
        \end{array} \right|
        \item \left|\begin{array}{cc}
            \log(a) &     \log(b) \\
            \sfrac{1}{2} & \sfrac{1}{4}
        \end{array} \right|
        \item \left|\begin{array}{cc}
            2m^2 & 2m^4-m \\
            m &  m^3-1
        \end{array} \right|
    \end{mathlist}
    \begin{answers}
        \begin{mathlist}[3]
            \item -10
            \item -12
            \item 1
            \item \log\frac{\sqrt[4]{a}}{\sqrt{b}}
            \item - m^{2}
        \end{mathlist}
    \end{answers}
\end{exercicio}

% 2
\begin{exercicio}
    Calcule os determinantes
    \begin{mathlist}[2]
        \item \left|\begin{array}{rrr}
            1 & 1 & 0 \\
            0 & 1 & 0 \\
            0 & 1 & 1
        \end{array} \right|
        \item \left|\begin{array}{rrr}
            1 & 3 &  2 \\
            -1 & 0 & -2 \\
            2 & 5 &  1
        \end{array} \right|
        \item \left|\begin{array}{rrr}
            -3 & 1 & -7 \\
            2 & 1 &  3 \\
            5 & 4 &  2
        \end{array} \right|
        \item \left|\begin{array}{rrr}
            9 & 7 & 11 \\
            -2 & 1 & 13 \\
            5 & 3 &  6
        \end{array} \right|
        \item \left|\begin{array}{rrr}
            0 & a & c \\
            -c & 0 & b \\
            a & b & 0
        \end{array} \right|
        \item \left|\begin{array}{rrr}
            2 & -1 & 0 \\
            m &  n & 2 \\
            3 &  5 & 4
        \end{array} \right|
    \end{mathlist}
  \begin{answers}
      \begin{mathlist}[3]
        \item 1
        \item -9
        \item 20
        \item 121
        \item b \left(a - c\right) \left(a + c\right)
        \item 4 m + 8 n - 26
    \end{mathlist}
  \end{answers}
\end{exercicio}


% 3
\begin{exercicio}
    Calcule os determinantes das matrizes usando o Teorema de Laplace
    \begin{mathlist}
        \item M = \left|\begin{array}{rrrr}
            3 & 4 &  2 &  1 \\
            5 & 0 & -1 & -2 \\
            0 & 0 &  4 &  0 \\
            -1 & 0 &  3 &  3
        \end{array} \right|
        \item M = \left|\begin{array}{rrrr}
            0 & a & b & 1 \\
            0 & 1 & 0 & 0 \\
            a & a & 0 & b \\
            1 & b & a & 0
        \end{array} \right|
        \item M = \left|\begin{array}{rrrr}
            1 & 3 & 2 & 0 \\
            3 & 1 & 0 & 2 \\
            2 & 3 & 0 & 1 \\
            0 & 2 & 1 & 3
        \end{array} \right|
        \item M = \left|\begin{array}{rrrrr}
            1 & 2 &  3 & 4 & 5 \\
            0 & a & -1 & 3 & 1 \\
            0 & 0 &  b & 2 & 3 \\
            0 & 0 &  0 & c & 2 \\
            0 & 0 &  0 & 0 & d
        \end{array} \right|
        \item M = \left|\begin{array}{cccccc}
            x & 0 & 0 & 0 & 0 & 0 \\ 
            a & y & 0 & 0 & 0 & 0 \\ 
            l & p & z & 0 & 0 & 0 \\ 
            m & n & p & x & 0 & 0 \\ 
            b & c & d & e & y & 0 \\ 
            a & b & c & d & e & z
        \end{array} \right|
    \end{mathlist}        
  \begin{answers}
      \begin{mathlist}[3]
        \item -208
        \item a^{2} + b^{2}
        \item 48
        \item a b c d
        \item x^{2} y^{2} z^{2}
    \end{mathlist}
  \end{answers}
\end{exercicio}

\begin{exercicio}
    Calcule o valor do determinante
    \( \left|\begin{array}{cccc}
    	1 & a & a & a \\
    	1 & 0 & a & a \\
    	1 & 1 & 0 & a \\
    	1 & 1 & 1 & 1
    \end{array} \right| \)
    \begin{answers}
        \(\qquad a^{2} \left( 1 - a \right) \)
    \end{answers}
\end{exercicio}

\begin{exercicio}
    Calcule o valor de $A$ sabendo que
    \[A = \left|\begin{array}{cccc}
    	1 & 2 & 4  &  8  \\
    	1 & 3 & 9  & 27  \\
    	1 & 4 & 16 & 64  \\
    	1 & 5 & 25 & 125
    \end{array} \right|\]
    \begin{answers}
          \(\qquad A = 12\)
    \end{answers}

\end{exercicio}

    
    \begin{exercicio}
        Determine $x$ tal que
        \begin{mathlist}
            \item \left|\begin{array}{cc}
            	2x & 3x+2 \\
            	 1 &    x
            \end{array} \right| = 0
            \item \left|\begin{array}{cc}
            	 2x  & x-2  \\
            	4x+5 & 3x-1
            \end{array} \right| = 11
        \end{mathlist}
        \begin{answers}
            \begin{mathlist}[2]
                \item x_1 = - \sfrac{1}{2} \quad x_2 = 2
                \item x_1 = -1 \quad x_2 = \sfrac{1}{2}
            \end{mathlist}
        \end{answers}
    \end{exercicio}
    
    \begin{exercicio}
        Determine $x$ tal que
        \begin{mathlist}
            \item \left|\begin{array}{rcr}
            	1 &  x  & x \\
            	2 & 2x  & 1 \\
            	3 & x+1 & 1
            \end{array} \right| = 0
            \item \left|\begin{array}{rrr}
            	1 &  x & 1 \\
            	1 & -1 & x \\
            	1 & -x & 1
            \end{array} \right| = 0
            \item \left|\begin{array}{rrr}
            	 1 &  x &  2 \\
            	-2 &  x & -4 \\
            	 1 & -3 & -x
            \end{array} \right| = 0
            \item \left|\begin{array}{ccr}
            	x-1 &  2  &  x \\
            	 0  &  1  & -1 \\
            	3x  & x+1 & 2x
            \end{array} \right| = 
            \left|\begin{array}{cc}
            	3x & 2x \\
            	4  & -x
            \end{array} \right|
        \end{mathlist}
        \begin{answers}
            \begin{mathlist}[2]
                \item x_1 = \sfrac{1}{2}
                \item x_1 = 0 \quad x_2 = 1
                \item x_1 = -2 \quad x_2 = 0
                \item x_1 = - \sfrac{\sqrt{3}}{3} \quad x_2 = \sfrac{\sqrt{3}}{3}
            \end{mathlist}
        \end{answers}
    \end{exercicio}

    \begin{exercicio}
        Encontre os valores de $a$ para os quais a matriz $A$ possui inversa.
        \[A = \left[ 
            \begin{array}{rrr}
            	1 & 1 & 0 \\
            	1 & 0 & 0 \\
            	1 & 2 & a
            \end{array}
        \right] \]
        \begin{answers} \(\qquad a\neq 0 \)
        \end{answers}
    \end{exercicio}
    
    \begin{exercicio}
        Calcule o determinante da matriz $M$ e determine para quais valores de $x$ ela possui inversa
        \[M = \left[ 
            \begin{array}{cccc}
            	4 & x & 0 & x \\
            	3 & 0 & 2 & 0 \\
            	0 & 2 & 0 & 0 \\
            	x & 0 & 1 & 0
            \end{array}
        \right]\]
        \begin{answers} \(\qquad x\neq 0 \) e \( x\neq \sfrac{3}{2} \)
        \end{answers}
    \end{exercicio}

    \begin{exercicio}
        Calcule o determinante da matriz 
        \[ A = \left[\begin{array}{cc}
            	e^{rt}  &   te^{rt}    \\
            	re^{rt} & (1+rt)e^{rt}
            \end{array} \right] 
        \]
        \begin{answers} \(\qquad e^{2 r t}\)
        \end{answers}
    \end{exercicio}
    
    \begin{exercicio}
        Sabendo que $\det A = -3$ calcule
        \begin{mathlist}[3]
            \item \det A^2
            \item \det A^3
            \item \det A^t
        \end{mathlist}
        \begin{answers}
            \begin{mathlist}[3]
                \item 9
                \item -27
                \item(-3)^t
            \end{mathlist}
        \end{answers}
    \end{exercicio}
     
     \begin{exercicio}
         Determine os valores de $\lambda$ tais que \[\det(A-\lambda I) = 0\] para as matrizes
         \begin{mathlist}
            \item A = \left[
                          \begin{array}{ccc}
                              0 & 1 & 2 \\ 
                              0 & 0 & 3 \\ 
                              0 & 0 & 0
                          \end{array}   
                      \right]
            \item A = \left[
                        \begin{array}{rrr}
                             1 & 0 &  0 \\ 
                            -1 & 3 &  0 \\ 
                             3 & 2 & -2
                        \end{array}   
                      \right]
        \end{mathlist}
    \end{exercicio}
     
    \begin{exercicio}
        Calcule o determinantes das matrizes
        \begin{mathlist}[2]
            \item A = \left[
                        \begin{array}{cc}
                            4 & 2 \\ 
                            3 & 2
                        \end{array} 
                      \right]
            \item A = \left[
                        \begin{array}{cc}
                          1 & 2 \\ 
                          0 & 4
                        \end{array} 
                      \right]
            \item AB
            \item 2A
            \item A^tB
            \item A^2B^{-1}
        \end{mathlist}
    \end{exercicio}

    \begin{exercicio}
        Calcule o determinante da matriz 
        \( \left[\begin{array}{cc}
        	a & b \\
        	b & a
        \end{array} \right] \), 
        sabendo que $2a = e^x + e^{-x}$ e $2b = e^x - e ^{-x}$.
    \end{exercicio}
    
    \begin{exercicio}
        Seja \( A = \left[\begin{array}{ccc}
        	2 & n & 2 \\
        	5 & 1 & m \\
        	7 & 3 & 7
        \end{array} \right] \)
        \begin{alphalist}
            \item Ache um valor de $m$, tal que $\det A = 0$, qualquer que seja $n$.
            \item Esse valor de $m$ é o único com essa propriedade?
        \end{alphalist}
    \end{exercicio}
    
    \begin{exercicio}
        A matriz real $A = (a_{ij})$ quadrada e de ordem~3, é definida pela regra
        \[a_{ij} = \begin{cases}
        	0,            & \text{se } i > j \\
        	i^j,          & \text{se } i = j \\
        	(-1)^{i + 1}, & \text{se } i < j
        \end{cases}\]
        \begin{alphalist}
            \item Escreva a matriz $A$ na forma usual de linhas e colunas, onde o valor numérico de $a_{ij}$ está na $i$-ésima linha e na $j$-ésima coluna de $A$.
            \item Calcule o determinante $D$ da matriz $A$.
        \end{alphalist}
    \end{exercicio}
    
    \begin{exercicio}
        Determine o valor de $x$ para que o determinante da matriz $C = AB^t$ seja igual a $602$, onde
        \[A = \left[\begin{array}{rrr}
        	1 & 2 & -3 \\
        	4 & 1 &  2
        \end{array} \right]
        \qquad
        B = \left[\begin{array}{crr}
        	x - 1 & 8 & -5 \\
        	 -2   & 7 &  4
        \end{array} \right]\]
        e $B^t$ é a matriz transposta de $B$.
    \end{exercicio}
    
    \begin{exercicio}
        Duas matrizes de ordem $n$ são inversas quando o produto delas é a identidade de ordem $n$. A inversa da matriz
        \[M = \left[\begin{array}{rrr}
        	1 &  1 & 2 \\
        	1 & -1 & 3 \\
        	2 &  3 & 1
        \end{array} \right]\]
        é a matriz
        \[
        M^{-1} = \left[\begin{array}{ccc}
        	x & 1 & 1 \\
        	1 & y & z \\
        	1 & z & w
        \end{array} \right]\]
        Encontre
        \begin{alphalist}
            \item o valor da expressão $A = x - 3y - z + w$.
            \item $\det(2M)$.
        \end{alphalist}
    \end{exercicio}
    
    \begin{exercicio}
        Determine o conjunto solução da equação
        \[\left|\begin{array}{lrc}
        	3^x &  1 & 0 \\
        	4   &  2 & 2 \\
        	2   & -1 & 3
        \end{array} \right| = \left|\begin{array}{cc}
            4 & 2 \\ 
            2 & 5
        \end{array} \right|\]
        \begin{answers} \( \qquad x = 1\)
        \end{answers}
    \end{exercicio}
    
    \begin{exercicio}
        Determine $x$ de modo que
        \[\left|\begin{array}{crl}
        	1 &  1 & 1   \\
        	2 & -3 & x   \\
        	4 &  9 & x^2
        \end{array} \right| \geq 0 \]
    \end{exercicio}
    
    \begin{exercicio}
        Resolva a equação
        \[\left|\begin{array}{cccc}
            1 & 1 & 0 & 1 \\ 
            x & 1 & x & 0 \\ 
            x & x & 1 & 0 \\ 
            0 & 1 & 0 & 1
        \end{array} \right| = 0\]
        \begin{answers} \(\quad x_1 = -1 \quad x_2 = 1 \)
        \end{answers}
    \end{exercicio}
    
    \begin{exercicio}
        Desenvolva o determinante a seguir
        \[x = \left|\begin{array}{rrrr}
        	a & b  &  c & d \\
        	0 & 1  &  2 & 0 \\
        	1 & -1 &  0 & 1 \\
        	0 & 1  & -1 & 0
        \end{array} \right|\]
    \end{exercicio}
    
    \begin{exercicio}
        Sabendo que
        \( \left|\begin{array}{ccc}
        	a & x & 1 \\
        	b & y & 2 \\
        	c & z & 3
        \end{array} \right| = 4  \) \\
        calcule
        \( M = \left|\begin{array}{ccc}
        	2a & 3x & -1 \\
        	2c & 3z & -3 \\
        	2b & 3y & -2
        \end{array} \right|  \)
    \end{exercicio}
    
    \begin{exercicio}
        Sem calcular os determinantes a seguir, justifique a sua nulidade
        \begin{mathlist}[2]
            \item \left|\begin{array}{rrr}
            	2 & 0 & -1 \\
            	5 & 0 &  5 \\
            	4 & 0 &  2
            \end{array} \right|
            \item \left|\begin{array}{rrr}
            	 a &  b &  c \\
            	ka & kb & kc \\
            	 x &  y &  m
            \end{array} \right|
            \item \left|\begin{array}{rrr}
            	 5 &  4 & 3 \\
            	-4 & -2 & 2 \\
            	 4 &  5 & 6
            \end{array} \right|
            \item \left|\begin{array}{rrr}
            	1 & 2 & 5 \\
            	m & 5 & 4 \\
            	1 & 2 & 5
            \end{array} \right|
            \item \left|\begin{array}{rrr}
            	2 & 3 &  4 \\
            	1 & 2 &  7 \\
            	3 & 5 & 11
            \end{array} \right|
            \item \left|\begin{array}{rrrr}
            	 1 &  2 &  3 &  4 \\
            	 5 &  6 &  7 &  8 \\
            	 9 & 10 & 11 & 12 \\
            	13 & 14 & 15 & 16
            \end{array} \right|
        \end{mathlist}
    \end{exercicio}
    
    \begin{exercicio}
        Calcule $x$ para que os determinantes das matrizes sejam iguais 
        \[A=\left[\begin{array}{crcc}
        	x + 1 &  0 & 0 & 2 \\
        	  0   & -2 & 0 & 5 \\
        	  x   &  0 & 3 & 0 \\
        	  0   &  1 & 0 & 4
        \end{array} \right] 
        \qquad
        B=
        \left[\begin{array}{cc}
        	1  & x \\
        	78 & 0
        \end{array} \right]\]
    \end{exercicio}
    
    \begin{exercicio}
        Determine as condições que $x$ deve satisfazer para que a matriz $A$ seja inversível
        \[A = \left[\begin{array}{cccc}
            1 & 2 & 3 & 4 \\ 
            1 & 3 & x & 5 \\ 
            1 & 3 & 4 & 3 \\ 
            1 & 6 & 5 & x
        \end{array} \right]\]
    \end{exercicio}
    
    \begin{exercicio}
        Seja a matriz 
        \[A = \left[\begin{array}{ccc}
            k & k & k \\ 
            k & k & 5 \\ 
            k & 5 & 5
        \end{array} \right]\]
        Determine $k \in \R$ de modo que a matriz $A$ não admita inversa.
    \end{exercicio}
    
    \begin{exercicio}
        Dadas as matrizes
        \[A = \left[\begin{array}{cc}
        	1 & 2 \\
        	2 & 0
        \end{array} \right]
        \qquad
        B = \left[\begin{array}{cr}
        	4 & -4 \\
        	1 &  4
        \end{array} \right]\]
        calcule o determinante da matriz $(AB)^{-1}$.
    \end{exercicio}
     
\end{exercicios} % Determinante de Uma Matriz e Suas Propriedades

%------------------------------------------------------------------------------%

%------------------------------------------------------------------------------%
% \chapter{Matrizes}
%------------------------------------------------------------------------------%

%------------------------------------------------------------------------------%
\begin{exercicios}{Revisão}

    \begin{exercicio}
        Resolva a equação  matricial $AX = B$, dadas as matrizes
        \[A = \left[\begin{array}{rr}
        2 & -1 \\
        -3 &  2
        \end{array} \right]	
        \qquad
        B = \left[\begin{array}{rr}
        3 &  2 \\
        -6 & -5
        \end{array} \right]	\]
    \end{exercicio}
    
\end{exercicios}

%------------------------------------------------------------------------------%
\begin{exercicios}{Desafio}
    
    \begin{exercicio}
        Determinar a inversa de cada uma das matrizes
        \begin{mathlist}[2]
        \item A = \left[\begin{array}{rr}
            5 & 6 \\ 
            4 & 5
        \end{array} \right]
        \item B = \left[\begin{array}{rr}
            2 & 5 \\ 
            1 & 3
        \end{array} \right]
        \item C = \left[\begin{array}{rr}
            1 & 0 \\ 
            0 & 2
        \end{array} \right]
        \item D = \left[\begin{array}{rr}
            1 & -1  \\ 
            1 &  1
        \end{array} \right]
        \end{mathlist}
    \end{exercicio}
    
    \begin{exercicio}
        Calcule a inversa da matriz
        \( A = \left[\begin{array}{cc}
        2 & 5 \\
        1 & 3
        \end{array} \right] \)
    \end{exercicio}
    
    \begin{exercicio}
        Um aluno, ao inverter a matriz
        \[A = \left[\begin{array}{ccc}
        1 & a & b \\
        0 & c & d \\
        4 & e & f
        \end{array} \right]
        =\left[a_{ij}\right], \qquad 1 \leq	i, j \leq 3
        \]
        cometeu um engano, e considerou o elemento $a_{31}$ igual a $3$, de forma que acabou invertendo a matriz
        \[B = \left[\begin{array}{ccc}
        1 & a & b \\
        0 & c & d \\
        3 & e & f
        \end{array} \right]	= \left[b_{ij} \right] 
        \]
        Com esse engano, o aluno encontrou 
        \[B^{-1} = \left[\begin{array}{rrr}
        \frac{5}{2} & 0 & -\frac{1}{2} \\[2mm]
        3 & 1 &           -1 \\[2mm]
        -\frac{3}{2} & 0 &  \frac{1}{2}
        \end{array} \right]\]
        Determinar $A^{-1}$.
    \end{exercicio}
    
    \begin{exercicio}
        Considere o sistema
        \[\begin{cases}
        x - my       &= 1 - m \\[2mm]
        (1 + m)x + y &= 1
        \end{cases}\]
        \begin{alphalist}
            \item Prove que o sistema admite solução única para cada número real $m$.
            \item Determine $m$ para que o valor de $x$ seja o maior possível.
        \end{alphalist}
    \end{exercicio}
    
    \begin{exercicio}
        Dada a matriz 
        \[A = \left[\begin{array}{cc}
        -2 & 3 \\
        -1 & 2
        \end{array} \right]\]
        calcule a sua inversa $A^{-1}$. 
        A relação especial, que você deve ter observado ente $A$ e $A^{-1}$ acima, 
        seria também encontrada se calculássemos as matrizes inversas de
        \begin{mathlist}
            \item \left[\begin{array}{rr}
                -3 & 4 \\ 
                -2 & 3
            \end{array} \right]
            \item \left[\begin{array}{rr}
                -5 & 6 \\ 
                -4 & 5
            \end{array} \right]
            \item \left[\begin{array}{rr}
                -1 & 2 \\
                0 & 1
            \end{array} \right]
        \end{mathlist}
        Generalize e demonstre o resultado observado.
    \end{exercicio}
    
    \begin{exercicio}
        Calcule a matriz inversa da matriz
        \[M = \left[\begin{array}{rrr}
        1 & -1 & 0 \\
        2 &  3 & 0 \\
        0 &  1 & 4
        \end{array} \right]\]
    \end{exercicio}
    
\end{exercicios}

%------------------------------------------------------------------------------%


\writeAnswers{03-Matrizes/03-Matrizes-Exercicios-ans}

%------------------------------------------------------------------------------%
